\documentclass[11pt]{article}
\usepackage[T1]{fontenc}
\usepackage[margin=1in]{geometry}
\usepackage{amsmath}
\usepackage{amssymb}
\usepackage{xurl}
\usepackage[hidelinks]{hyperref}

\newcommand{\pr}{\operatorname{pr}}

\setlength{\emergencystretch}{2em}
\newenvironment{notes}
  {\begin{itemize}\setlength{\itemsep}{0.45em}\setlength{\parsep}{0pt}}
  {\end{itemize}}

\title{CorrPFN: Method and Experiment Notes}
\author{}
\date{}

\begin{document}
\maketitle

% Sections 1--8 of the agreed report structure.
% Add the later sections below.

\section{The EPIT Dataset and Its Splits}
\label{sec:epit-dataset}

\subsection{Source, Prediction Task, and Usable Rows}
\label{sec:dataset-source}

\begin{notes}
    \item \textbf{Source.} The data come from
    \textit{Electrochemical Metrics for Corrosion Resistant Alloys},
    released on
    \href{https://figshare.com/articles/dataset/Electrochemical_Metrics_for_Corrosion_Resistant_Alloys/13038257}{Figshare}
    and described by Nyby et al.\ in
    \href{https://doi.org/10.1038/s41597-021-00840-y}{\textit{Scientific Data} (2021)}.

    \item \textbf{Selected table and target.} This project uses the workbook
    table \texttt{Pitting Potential}. The value to predict is
    \texttt{Epit, mV (SCE) Avg.}: the reported average pitting potential.
    Here, mV means millivolts, and SCE means saturated calomel electrode,
    the reference electrode relative to which the potential is expressed.

    \item \textbf{Meaning of a row.} Each row links an alloy composition
    and reported test conditions to an EPIT value. The same or a similar
    composition can appear in several rows, for example when measurements
    use different environments or test methods.

    \item \textbf{Usable observations.} The source table has \(810\) rows.
    Of these, \(760\) contain a usable numeric value for the selected target.

    \item \textbf{Excluded observations.} The other \(50\) rows cannot be
    used as numeric targets: \(41\) contain the text
    \texttt{Transpassive}, and \(9\) contain \texttt{NA}.
    These entries are excluded rather than replaced with invented EPIT values.
\end{notes}

\subsection{What EPIT Means for Corrosion}
\label{sec:epit-meaning}

\begin{notes}
    \item \textbf{Pitting corrosion.} A protective surface film, called the
    \emph{passive film}, can break down locally. Corrosion then develops
    at these localized sites, forming pits.

    \item \textbf{Pitting potential.} EPIT is the reported electrochemical
    potential at which passive-film breakdown or sustained pitting is
    identified during a test. It describes a breakdown threshold, not a
    corrosion rate.

    \item \textbf{Meaning of a higher EPIT.} Under comparable conditions,
    a more positive EPIT generally means that a higher applied potential
    is needed to initiate sustained pitting. It therefore indicates
    greater resistance to pit initiation under those conditions.

    \item \textbf{Why chemistry and environment matter.} EPIT depends on
    the alloy composition and material state, together with temperature,
    chloride concentration, pH, and other solution conditions. Surface
    preparation and previous exposure can also affect the measurement.

    \item \textbf{Why the test method matters.} Different procedures can
    report breakdown at different potentials. The workbook combines
    potentiodynamic, potentiostatic, scratch-based, and other procedures.
    Scan rate, stabilization time, current thresholds, and the criterion
    used to identify breakdown can differ between studies.

    \item \textbf{How to compare values.} A higher recorded EPIT does not,
    by itself, prove that one alloy is intrinsically more resistant than
    another. The environments, reference electrodes, surface conditions,
    and test procedures must also be considered.

    \item \textbf{What the inputs cannot fully describe.} The retained
    columns do not separately record every relevant detail of
    microstructure, surface preparation, exposure history, or test
    procedure. Some differences between measurements therefore remain
    unexplained by the available inputs.
\end{notes}

\subsection{Input Features and Alloy Classes}
\label{sec:dataset-features}

\begin{notes}
    \item \textbf{Input schema.} Each row has \(21\) retained input columns:
    \[
        21
        = 17\text{ composition}
        + 3\text{ environment}
        + 1\text{ test method}.
    \]
    The reported models receive the same raw input columns.

    \item \textbf{Composition columns.} The \(17\) elements are kept in
    this fixed order:
    \[
    \begin{gathered}
        \mathrm{Fe,\ Cr,\ Ni,\ Mo,\ W,\ Nb,\ Al,\ V,\ Ta,}\\
        \mathrm{Re,\ Ce,\ Ti,\ Co,\ B,\ Mg,\ Y,\ Gd}.
    \end{gathered}
    \]
    Each element has its own column. For example, the model receives the
    chromium amount and the molybdenum amount separately; they are not
    combined into a single alloy score.

    \item \textbf{Composition units.} Element amounts are reported in
    weight percent, written wt.\%. For example, \(20\)~wt.\% chromium
    means \(20\)~g of chromium per \(100\)~g of alloy.

    \item \textbf{Environmental columns.}
    \begin{notes}
        \item \texttt{Test Temp. oC}: the test temperature in degrees
        Celsius.
        \item \texttt{[Cl-] M}: the chloride-ion concentration in the
        solution. The unit M means moles per litre.
        \item \texttt{[Cl-] pH}: the reported pH of the chloride-containing
        solution, describing its acidity or alkalinity.
    \end{notes}

    \item \textbf{Test-method column.}
    \texttt{[Cl-] Test Method} records the experimental method or
    protocol. It is classified as process/history information because
    the procedure can influence the measured EPIT. It is one recorded
    category, not a complete record of every procedural detail.

    \item \textbf{How columns were selected.} The evaluator filters
    columns by feature group, availability of usable values, and the
    number of categories in categorical columns. Availability of usable
    values is called \emph{coverage}; the number of distinct categories
    is called \emph{categorical cardinality}.

    \item \textbf{Omitted composition columns.} N, C, Si, Mn, Cu, P,
    and S fall below the required \(80\%\) finite-value coverage.
    They are therefore omitted from the model inputs. This is a
    data-availability decision; these elements can still be relevant
    to corrosion.

    \item \textbf{Alloy classes.} The dataset includes Fe alloys,
    Ni--Cr--Mo alloys, Al alloys, high-entropy alloys, and other alloys.
    These are material classes, separate from the composition groups
    constructed below.
\end{notes}

\subsection{Building Composition Groups}
\label{sec:composition-groups}

\begin{notes}
    \item \textbf{Reason for grouping.} Several rows can describe the
    same or very similar alloy compositions. These rows are kept
    together so that a composition and a closely related version do
    not end up on opposite sides of a split.

    \item \textbf{Values used for grouping.} Grouping uses the \(17\)
    retained composition columns. For this operation only, missing
    values are treated as zero and each value is rounded to
    \(0.01\)~wt.\%.

    \item \textbf{No rescaling to \(100\%\).} The retained composition
    values are not renormalized before grouping. The seven omitted
    elements can account for part of the alloy, so the \(17\) visible
    amounts need not sum to \(100\%\).

    \item \textbf{Composition distance.} Compare two rows by adding
    the absolute differences in their element amounts:
    \begin{equation}
        d(i,j)=\sum_{e=1}^{17}|x_{i,e}-x_{j,e}|.
        \label{eq:composition-distance}
    \end{equation}
    Here, \(x_{i,e}\) is the rounded weight-percent value of element
    \(e\) in row \(i\). This is called the \emph{L1 distance}.
    For example, differences of \(0.6\)~wt.\% in one element and
    \(0.3\)~wt.\% in another, with all other retained amounts equal,
    give \(d=0.9\)~wt.\%.

    \item \textbf{Linking rule.} Two rows are linked when
    \[
        d(i,j)\leq1.0\ \text{wt.\%}.
    \]
    All rows connected by these links form one group. The technical
    name for such a group is a \emph{connected component}.

    \item \textbf{Why a group can span more than the threshold.}
    If \(A\) is close to \(B\), and \(B\) is close to \(C\), all three
    stay together even if \(A\) and \(C\) are more than \(1.0\)~wt.\%
    apart. This ensures that every directly linked pair stays on the
    same side of every split.

    \item \textbf{Resulting groups.} The frozen split has \(301\)
    composition groups. The largest group contains \(45\) rows.
    The largest within-group composition span is \(5.0\)~wt.\%;
    this is possible because of the chain of links described above.

    \item \textbf{Meaning of composition separation.} The split tests
    prediction for composition groups absent from the labeled examples
    supplied as context. A labeled example contains both the input
    features and a known EPIT value.

    \item \textbf{Scope of that separation.} These groups are separated
    between the labeled partitions. Their compositions are not
    necessarily absent from every source used by the project:
    the empirical feature profiles used for synthetic generation also
    include final-test features.
\end{notes}

\subsection{Development Partition, Final-Test Partition, and Folds}
\label{sec:dataset-partitions}

\begin{notes}
    \item \textbf{Outer split.} Whole composition groups are assigned
    to two partitions:
    \[
        760\text{ rows}
        =
        \underbrace{608\text{ development rows}}_{80\%}
        +
        \underbrace{152\text{ final-test rows}}_{20\%}.
    \]
    Development rows support method development. The final-test
    partition is used to assess the selected model.

    \item \textbf{Five development folds.} The development partition
    is divided into five fixed subsets called \emph{folds}. Each fold
    contains \(121\)--\(122\) rows.

    \item \textbf{Role of a fold.} For one development evaluation,
    one fold supplies the validation rows to predict. The other four
    folds supply the \(486\)--\(487\) labeled context rows.

    \item \textbf{Example.} First, fold~1 can be the validation subset,
    with folds~2--5 as context. Next, fold~2 can be the validation
    subset, with the other four as context. Repeating this for all five
    folds gives every development row one turn as a validation row.

    \item \textbf{Group boundaries are respected throughout.}
    A composition group cannot be divided between development and
    final test, or between development folds. No pair linked by the
    distance rule crosses either type of boundary.

    \item \textbf{Balancing the partitions.} With the row counts fixed,
    the assignment aims to balance the following properties across the
    outer split and the development folds:
    \begin{notes}
        \item Alloy class and broad test-method family.
        \item EPIT values, divided into ten groups using target
        quantiles (\emph{deciles}).
        \item Composition isolation, divided into five quantile
        groups (\emph{quintiles}). Isolation is the distance from a
        composition group to its nearest other group.
        \item Temperature, chloride concentration, and pH.
        \item Source reference, so the source studies are represented
        across the partitions.
    \end{notes}

    \item \textbf{Environmental bins used for balancing.}
    Temperature uses boundaries at \(10\), \(35\), and
    \(70\,^{\circ}\mathrm{C}\); chloride concentration uses
    \(0.01\), \(0.1\), \(0.6\), and \(1.0\)~M, with a separate zero
    category; and pH uses \(3\), \(6\), and \(8\).
    Missing environmental values form separate categories.

    \item \textbf{Source-reference categories.} The ten most frequent
    source references are kept as separate categories. The remaining
    references are pooled into an \texttt{other} category.

    \item \textbf{Use of targets when constructing the split.}
    EPIT values from the usable corpus help balance the partitions.
    The split is therefore not \emph{target-blind}: it is not
    constructed using input features alone. This use of EPIT values
    concerns the assignment of rows to partitions.

    \item \textbf{Source overlap.} Measurements from one source study
    may appear in both development and final test, provided their
    composition groups remain separate. The split therefore does not
    test generalization to entirely new publications or laboratories.

    \item \textbf{Fixed assignments.} The chosen groups, row
    assignments, and fold identifiers are saved in frozen split
    artifacts. Later steps verify these records so that they all use
    the same split.
\end{notes}

\subsection{Preparing Real Inputs After Splitting}
\label{sec:real-data-representation}

\begin{notes}
    \item \textbf{Numeric inputs.} The \(17\) composition and three
    environmental columns are converted to numeric values.

    \item \textbf{Missing numeric values.} For TabICL evaluation,
    a missing value is replaced by the mean of that column calculated
    from the context/training rows. The same fitted means are applied
    to the validation or final-test rows.

    \item \textbf{Example of context-fitted imputation.} If the
    context rows have a mean test temperature of
    \(25\,^{\circ}\mathrm{C}\), a missing temperature is filled with
    \(25\,^{\circ}\mathrm{C}\). Temperatures in the validation or
    final-test rows do not change that fitted mean. This evaluation
    treatment differs from the zero replacement used only for
    composition grouping.

    \item \textbf{Encoding the test method.} TabICL receives one
    integer-coded method column. For illustration,
    \[
        \text{method A}\mapsto0,\qquad
        \text{method B}\mapsto1,\qquad
        \text{method C}\mapsto2.
    \]
    One-hot encoding would instead create a separate yes/no column
    for each method; it is not used here.

    \item \textbf{Categories available in one split.} The complete
    task contains \(52\) categories, including a missing-value category.
    The evaluation mapping is learned only from the context/training
    rows, so a particular mapping can contain fewer categories.
    For example, a context containing only \(40\) methods provides a
    mapping for those observed methods.

    \item \textbf{Unseen and missing methods.} A method found in a
    validation or final-test row but absent from the fitted mapping
    receives the reserved code \(-1\). The encoder also uses
    \(-1\) for missing values.

    \item \textbf{Meaning of the codes.} The integers are
    \emph{nominal identifiers}: they identify categories without
    ranking them. Method~2 is not physically higher than method~1,
    and adjacent codes do not mean that two methods are more similar.
    Treating these labels as numeric inputs remains an approximation.

    \item \textbf{Column order.} After preprocessing, the original
    \(21\)-column order is restored. Each element, environmental
    variable, and test-method column therefore stays in its designated
    position.
\end{notes}

\section{How TabICL Learns and Makes Predictions}
\label{sec:tabicl}

\subsection{The Prediction Task and the Presentation's Notation}
\label{sec:pfn-notation}

\begin{notes}
    \item \textbf{Sources and notation.} This section follows Philipp
    Benner's presentation \textit{Prior-Data Fitted Networks}
    (10 June 2026), stored in
    \path{resources_pres/21_-_pfn.pdf}. The accompanying notes in
    \path{resources_pres/Baysian PFN.docx} supply further explanations.
    References to slide numbers below use the numbered slides, not
    the PDF page numbers.

    \item \textbf{Observed examples.} Following the slides, write
    \[
        x=(x_i)_{i=1}^{n},
        \qquad
        y=(y_i)_{i=1}^{n}.
    \]
    The pair \((x,y)\) is the labeled context: \(x_i\) contains the
    inputs for one observed row, and \(y_i\) is its known response.
    For this project, \(y_i\) is EPIT.

    \item \textbf{The row to predict.} The new input is \(x^\ast\),
    called the \emph{query}. Its unknown response is \(y^\ast\).
    The star marks a query quantity; it does not mean multiplication.

    \item \textbf{Tabular notation.} When the context inputs are
    arranged as a matrix, the presentation writes
    \(X\in\mathbb{R}^{n\times p}\), with
    \(x^\ast\in\mathbb{R}^{p}\). Here, \(n\) is the number of context
    rows and \(p\) is the number of input columns. The reported
    CorrPFN model uses \(p=21\). Thus, \((x,y)\) and \((X,y)\)
    describe the same context in collection and matrix notation.

    \item \textbf{Task mechanism: \(\theta\).} In the PFN
    interpretation, \(\theta\) describes one possible mechanism
    that could generate a task: for example, its function, dependency
    structure, generator parameters, and noise. It is neither a
    dataset nor the weights of the TabICL transformer.

    \item \textbf{Transformer weights: \(\phi\).} These are the
    parameters learned during PFN pretraining. One set of weights
    \(\phi\) is shared across many sampled mechanisms \(\theta\).

    \item \textbf{Two predictive distributions.}
    \(\pr(y^\ast\mid x^\ast,x,y)\) denotes the predictive distribution
    implied by the statistical model and its prior.
    \(q_\phi(y^\ast\mid x^\ast,x,y)\) denotes the distribution
    produced by the trained PFN. The symbol \(\mid\) means
    \emph{given}. For a continuous response such as EPIT, these
    expressions can denote probability densities.

    \item \textbf{A distribution rather than just one value.}
    A predictive distribution describes which response values are
    plausible and how concentrated or spread out they are.
    A mean, median, or prediction interval can then be calculated
    from that distribution.
\end{notes}

\subsection{How This Differs from the Usual Machine-Learning Workflow}
\label{sec:traditional-ml-versus-pfn}

\begin{notes}
    \item \textbf{Usual fit-then-predict workflow.} For a supervised
    neural network or a model such as CatBoost, the labeled dataset
    is used to fit a model. In generic parameter notation,
    \[
        (x,y)\xrightarrow{\text{fit}}\widehat\theta,
        \qquad
        \widehat y^\ast=f_{\widehat\theta}(x^\ast).
    \]
    The data determine a fitted parameter setting
    \(\widehat\theta\), which is then used to predict new rows.

    \item \textbf{What happens for another dataset.} In that workflow,
    a new prediction task normally requires fitting or adapting a
    model to its labeled examples. The fitted model may itself be
    an ensemble, and it may output probabilities; one fitted model
    does not necessarily mean one tree or one point-valued output.

    \item \textbf{PFN pretraining.} A Prior-Data Fitted Network is
    trained in advance on many synthetic prediction tasks. Its
    weights \(\phi\) learn a procedure for using a labeled context
    to predict a query, rather than only a relationship fitted to
    one observed dataset.

    \item \textbf{PFN use on a new context.} The labeled examples and
    the query are supplied together as inputs:
    \[
        (x,y,x^\ast)
        \xrightarrow{\text{frozen PFN}}
        q_\phi(y^\ast\mid x^\ast,x,y).
    \]
    The prediction can change when the context changes, even though
    the transformer weights \(\phi\) remain fixed.

    \item \textbf{Where the adaptation happens.} Both workflows can
    use fixed weights when predicting a query. The difference is how
    they use the labeled examples for a new task: the conventional
    workflow fits model parameters, whereas the PFN uses those
    examples inside its forward computation. This is
    \emph{in-context learning}.

    \item \textbf{Scope of the comparison.} This describes the usual
    fitted-model workflow, not every traditional method. Methods
    such as nearest-neighbour prediction also use stored examples
    at prediction time. The distinctive PFN idea is to learn the
    context-to-prediction procedure through pretraining across tasks.

    \item \textbf{Computational trade-off.} PFNs invest computation
    in pretraining and reuse the resulting prediction procedure.
    They still have to process the context during inference.
    This does not imply that every PFN prediction is cheaper than
    a prediction from an already fitted conventional model.
\end{notes}

\subsection{Bayesian Interpretation: Prior, Likelihood, Posterior, and MAP}
\label{sec:bayesian-updating}

The previous subsection explained that a PFN uses labeled examples
as context. We now ask what prediction that context should imply
under a statistical model. This subsection and the next develop
the Bayesian reference calculation; they are not additional
inference steps executed by CorrPFN. Section~\ref{sec:pfn-pretraining}
then explains how synthetic training teaches a network to
approximate the resulting prediction.

\begin{notes}
    \item \textbf{Prior: \(\pr(\theta)\).} The prior describes which
    mechanisms are plausible before observing the current context.
    For a synthetic task generator, it determines which kinds of
    mechanisms are sampled and how often.

    \item \textbf{The prior in this project is a sampling recipe.}
    Instead of specifying only a simple density such as a normal
    prior on regression coefficients, CorrPFN specifies a program
    that samples task mechanisms and observations. For the selected
    Trial~56, the recipe includes these fixed settings, rounded here:
    \begin{notes}
        \item An informed task is chosen with probability
        \(\rho=0.75\); otherwise a generic task is chosen.
        \item Within informed tasks, a multilayer perceptron (MLP)
        and a tree are equally likely as target generators. The
        generic branch uses
        MLP/tree probabilities of \(0.7/0.3\).
        \item Informed rows use empirical composition templates
        and environmental/test-method records. Positive element
        amounts are multiplied by \(\exp(\epsilon)\), where
        \(\epsilon\sim\operatorname{Normal}(0,0.1028^2)\),
        before full-composition normalization.
        \item Each informed response combines a standardized
        random-generator response with a standardized analytical
        response, using coefficients \(0.7952\) and \(0.2048\),
        respectively, and then standardizes the result again.
    \end{notes}
    These are actual selected settings, not illustrative choices.
    The complete recipe appears in
    Section~\ref{sec:informed-prior}. Its response coefficients
    describe a weighted combination, not shares of explained variance.

    \item \textbf{A fixed recipe still generates different mechanisms.}
    One possible draw selects the informed branch, a particular
    random MLP, and the linear PREN-inspired analytical rule.
    Another can select a tree and a different rule. The first
    rule is sampled with probability about \(0.1342\); its
    saved weights for material protection, environmental
    aggressiveness, and material--chloride interaction are
    approximately \((0.5560,0.4013,0.0427)\).
    The rule selection can change between tasks, while that
    rule's coefficient vector stays fixed. The selected response
    mixture and perturbation strength also stay fixed during
    the run; Optuna chose them beforehand.

    \item \textbf{Keep the recipe, mechanism, data, and predictor separate.}
    The recipe defines which mechanisms can be drawn and how often.
    A particular draw of the task-level choices, random function,
    and noise settings belongs to \(\theta\). The generated rows
    form a dataset \(D\), with observations varying under that
    mechanism. The transformer weights \(\phi\) are learned across
    many such datasets. If the generator is an MLP, its random
    weights belong to the mechanism; they are different from the
    trained weights of the prediction transformer. The seven
    analytical rules are therefore only one part of the space of
    possible mechanisms.

    \item \textbf{An example of an observation drawn from the recipe.}
    The saved composition bank includes an Fe-alloy template with
    \(16\) wt.\% Cr, \(11\) wt.\% Ni, and \(3\) wt.\% Mo.
    An illustrative chromium perturbation of \(\epsilon=0.05\)
    changes its amount to \(16e^{0.05}\approx16.82\) wt.\%
    before normalization; other positive amounts receive their own
    perturbations. A sampled environment and method complete the
    row. Its response is generated using the task's random function
    and selected analytical rule. It is a standardized synthetic
    score; the original experimental EPIT label is not copied.

    \item \textbf{What the prior means after project development.}
    The recipe is fixed before conditioning for a prediction,
    but it has already been adapted using empirical feature
    profiles, development-label rule calibration, and development
    selection. Thus, the word prior does not imply that no real
    data influenced its construction. The relevant data boundaries
    are documented in Sections~\ref{sec:feature-profile-provenance}
    and~\ref{sec:development-score-interpretation}.

    \item \textbf{Likelihood: \(\pr(x,y\mid\theta)\).} The likelihood
    measures how well a proposed mechanism explains the observed
    context. The observations are held fixed while different
    \(\theta\) values are compared. It is not itself a probability
    distribution over \(\theta\).

    \item \textbf{Posterior: \(\pr(\theta\mid x,y)\).}
    Bayes' theorem combines the prior and likelihood:
    \begin{equation}
        \pr(\theta\mid x,y)
        =
        \frac{\pr(x,y\mid\theta)\pr(\theta)}
             {\pr(x,y)}.
        \label{eq:bayes-posterior}
    \end{equation}
    In words: the plausibility after seeing the context is
    proportional to the prior plausibility multiplied by how well
    the mechanism explains the context. The denominator normalizes
    the result.

    \item \textbf{A corrosion interpretation.} Imagine candidate
    mechanisms with different strengths of composition effects,
    environmental effects, and their interactions. Context
    observations can make some candidates more plausible than others.
    This is an illustration of statistical reasoning, not a claim
    that the real physical mechanism has been identified.

    \item \textbf{What compatibility with the context would mean.}
    Suppose otherwise comparable context rows show lower EPIT at
    higher chloride concentration. A mechanism that reproduces
    this pattern would receive more support than one that
    systematically contradicts it, taking its noise assumptions
    into account. This is about explaining the observed responses
    at the observed inputs, not about matching synthetic row
    identifiers or finding the nearest saved synthetic table.
    CorrPFN does not explicitly evaluate this likelihood or
    construct \(\pr(\theta\mid x,y)\) at inference. Its generator
    supplies samples from an implicit probability model even
    though evaluating that model's likelihood may be difficult.

    \item \textbf{MAP chooses one setting.} The
    \emph{maximum a posteriori} estimate is
    \begin{equation}
        \widehat\theta
        =
        \underset{\theta}{\arg\max}\,
        \pr(\theta\mid x,y)
        =
        \underset{\theta}{\arg\max}\,
        \pr(x,y\mid\theta)\pr(\theta).
        \label{eq:map-estimate}
    \end{equation}
    The expression \(\arg\max\) means the value of \(\theta\)
    that makes this expression largest. For continuous
    parameters, MAP selects the highest posterior density.

    \item \textbf{MAP is a comparison here, not CorrPFN's selection step.}
    It illustrates committing to one explanation after seeing a
    context. Optuna selecting Trial~56 instead chooses the generator
    recipe, and selecting step~\(2{,}500\) chooses transformer
    weights. Neither operation computes the MAP mechanism for a
    prediction context. The next subsection explains the alternative
    of retaining several plausible mechanisms in the prediction.

    \item \textbf{Connection to conventional fitting.} Maximizing
    the posterior is equivalent to minimizing
    \[
        -\log\pr(x,y\mid\theta)-\log\pr(\theta).
    \]
    A loss plus regularization can therefore be interpreted as MAP
    when its terms correspond to a negative log-likelihood and a
    negative log-prior. With a constant prior over the relevant
    parameter space, MAP reduces to maximum likelihood.
    Ordinary machine-learning fitting is not automatically MAP.

    \item \textbf{Bayesian neural networks are a different use of
    random weights.} The presentation writes \(W\sim\pr(W)\)
    for a Bayesian neural network and averages predictions using
    \(\pr(W\mid x,y)\). For PFNs, the Bayesian averaging discussed
    next concerns the task mechanism \(\theta\).
    The inference network has its own weights \(\phi\), which
    are not sampled from a transformer-weight posterior at test time.
    These distinctions follow slides~1--3 and~15.
\end{notes}

\subsection{Posterior Prediction: Combining Plausible Mechanisms}
\label{sec:posterior-prediction}

Bayesian updating gives a distribution over mechanisms. We now
turn that distribution into a prediction for a query. Several
mechanisms can explain the context reasonably well but disagree
about the query; posterior prediction retains that disagreement
instead of discarding it when choosing one mechanism.

\begin{notes}
    \item \textbf{Prediction using one estimated mechanism.}
    A plug-in prediction uses a single setting, such as MAP:
    \[
        \pr(y^\ast\mid x^\ast,\widehat\theta).
    \]
    It can describe response noise, but it does not average over
    uncertainty about the mechanism \(\theta\).

    \item \textbf{Posterior predictive distribution.}
    Following slide~3, combine the predictions from possible
    mechanisms:
    \begin{equation}
        \pr(y^\ast\mid x^\ast,x,y)
        =
        \int
        \pr(y^\ast\mid x^\ast,\theta)
        \pr(\theta\mid x,y)\,d\theta.
        \label{eq:posterior-predictive}
    \end{equation}
    The first factor is the prediction under one mechanism.
    The second is its posterior plausibility. The integral adds
    their contributions across mechanisms.

    \item \textbf{Assumptions behind this form.} The displayed
    expression uses the slides' convention that the query input
    is treated as given for mechanism inference, and that
    \(\theta\) and \(x^\ast\) suffice to predict \(y^\ast\).
    In a joint model where the query's feature values also provide
    information about \(\theta\), the mechanism posterior must
    additionally condition on \(x^\ast\).

    \item \textbf{What is being averaged.} The integral averages
    \emph{predictive distributions}, not the values of the
    parameters \(\theta\). It therefore retains contributions from
    several plausible explanations instead of choosing only one.

    \item \textbf{Two sources of predictive spread.}
    A particular mechanism can allow noisy responses even when
    its parameters are known. Different plausible mechanisms can
    also predict different responses at the same input. Averaging
    their distributions retains both within-mechanism response
    variation and between-mechanism disagreement. The latter is
    why selecting one mechanism can miss uncertainty even if its
    own predictive distribution includes noise.

    \item \textbf{Small illustrative example.} The accompanying
    Word notes consider three mechanisms with posterior weights
    \(0.2\), \(0.1\), and \(0.7\). Suppose their predictive means
    for one query are \(200\), \(350\), and \(300\)~mV.
    The combined predictive mean is
    \[
        0.2(200)+0.1(350)+0.7(300)=285\ \text{mV}.
    \]
    Choosing only the most plausible mechanism would instead use
    the third mechanism. These numbers illustrate averaging;
    they are not results from the EPIT experiment.

    \item \textbf{These weights differ from the seven rule probabilities.}
    The illustrative \(0.2,0.1,0.7\) weights express plausibility
    after a context has been observed. CorrPFN's saved rule
    probabilities, such as \(0.1342\) for linear PREN, determine
    frequencies during synthetic generation. They come from
    normalized development correlations and are not likelihoods
    or posterior probabilities for the current context. CorrPFN
    does not produce an explicit updated weight for each of the
    seven rules at inference.

    \item \textbf{Mean and median are different summaries.}
    The \(285\)~mV calculation gives a mean, not the complete
    predictive distribution or necessarily its median. The
    reported CorrPFN model uses the median of its predicted
    distribution as its point prediction.
    In general, the median of a mixture must be obtained from
    the combined distribution; averaging the component medians
    does not give that median.

    \item \textbf{Numerical Bayesian approximation.} The integral
    can be difficult to calculate. Slide~4 describes approximating
    the posterior or drawing samples
    \(\theta^{(1)},\ldots,\theta^{(S)}\) from it, giving
    \[
        \pr(y^\ast\mid x^\ast,x,y)
        \approx
        \frac{1}{S}\sum_{s=1}^{S}
        \pr(y^\ast\mid x^\ast,\theta^{(s)}).
    \]
    Here, \(S\) is the number of samples. Equal weights are
    appropriate because the samples represent the
    \emph{posterior}. Drawing mechanisms from the prior and
    averaging them equally does not, by itself, incorporate the
    observed context.

    \item \textbf{Where the implementation differs from this calculation.}
    Explicit posterior sampling is a general numerical approach,
    not a step in the reported CorrPFN prediction path. CorrPFN
    learns a direct context-to-distribution function instead.
    The next subsection explains why generated observations are
    enough to train that function without first computing the
    posterior-predictive integral for every training example.
\end{notes}

\subsection{How a PFN Learns the Predictive Distribution}
\label{sec:pfn-pretraining}

The preceding subsections defined a statistical prediction that
may be expensive to calculate explicitly. We now move to the
PFN training procedure: use the generator to create prediction
problems, then learn a function that maps each context and query
to an appropriate predictive distribution.

\begin{notes}
    \item \textbf{Learning the prediction procedure.} A PFN is
    trained to produce
    \begin{equation}
        q_\phi(y^\ast\mid x^\ast,x,y)
        \approx
        \pr(y^\ast\mid x^\ast,x,y).
        \label{eq:pfn-predictive-approximation}
    \end{equation}
    It learns this mapping across many synthetic tasks.
    At inference, it does not explicitly enumerate mechanisms,
    calculate their posterior weights, or retrieve its original
    synthetic training datasets.

    \item \textbf{Sample a mechanism, then sample observations.}
    Slides~7--8 describe the generative model
    \[
        \pr(\theta,x,y,x^\ast,y^\ast)
        =
        \pr(\theta)\pr(x,y,x^\ast,y^\ast\mid\theta).
    \]
    First draw \(\theta\sim\pr(\theta)\). Then draw the context
    and query from that mechanism. One mechanism can generate many
    different datasets because the observations and noise also vary.

    \item \textbf{The prior predictive distribution.}
    Averaging out the unobserved mechanism gives
    \[
        \pr(x,y,x^\ast,y^\ast)
        =
        \int
        \pr(x,y,x^\ast,y^\ast\mid\theta)\pr(\theta)\,d\theta.
    \]
    This is the distribution of observable training tasks.
    Sampling a mechanism and then a task produces samples from
    this distribution without calculating the integral.

    \item \textbf{Prior predictive and posterior predictive describe
    different stages.} The prior predictive distribution describes
    complete possible datasets before a particular context is
    observed. Once a context and query inputs are given, the
    relevant distribution is the conditional distribution of the
    missing query responses. This conditional distribution
    averages over the remaining uncertainty about the mechanism;
    the posterior-predictive expression in the preceding subsection
    describes that averaging under its stated assumptions.
    The two distributions come from the same generative model.

    \item \textbf{What the network sees.} It receives the context
    inputs \(x\), their labels \(y\), and the query input \(x^\ast\).
    The sampled query label \(y^\ast\) is withheld from the network
    input and used to score its prediction.

    \item \textbf{Why context and query must share a task.}
    In one generated task, the labeled rows demonstrate the
    relationship that also produces its query responses. In
    another task, that relationship can change. The transformer
    must therefore learn how to use the context to predict under
    the current relationship. Merely learning one fixed mapping
    from alloy features to a target would not solve all these
    changing tasks.

    \item \textbf{The presentation's training objective.}
    For probabilistic prediction, the slides use negative
    log-likelihood:
    \begin{equation}
        \mathcal{L}(\phi)
        =
        \mathbb{E}
        \left[-\log q_\phi(y^\ast\mid x^\ast,x,y)\right].
        \label{eq:pfn-nll}
    \end{equation}
    The expectation is an average over sampled tasks and their
    context/query observations. Training adjusts \(\phi\) so that
    the generated query outcomes receive better predictions.

    \item \textbf{Why this can learn a Bayesian prediction.}
    For a fixed context and query, the expected log loss is
    minimized by the true conditional distribution of the sampled
    tasks, if the model can represent it. Training on enough
    tasks therefore targets that predictive distribution.
    Exact posterior calculations are not required as training
    labels: the sampled outcomes \(y^\ast\) provide supervision.
    This is the construction described by
    \href{https://arxiv.org/abs/2112.10510}{M\"uller et al.,
    \textit{Transformers Can Do Bayesian Inference}}.

    \item \textbf{Why individual outcomes can teach a whole distribution.}
    At a given context and query, a distributional scoring rule
    rewards assigning suitable probabilities to the outcomes that
    can occur. Across generated tasks, the network learns shared
    patterns of how those probabilities depend on the context.
    It need not encounter exactly the same continuous-valued table
    repeatedly. The training examples contain generated outcomes,
    not exact Bayesian predictive distributions supplied by a
    separate calculation. The local loss is CRPS rather than the
    slides' log loss; Section~\ref{sec:tabicl-regression} explains
    why it still targets a conditional distribution.

    \item \textbf{Several queries in one synthetic dataset.}
    Slide~32 writes
    \[
        D=\{(x_i,y_i)\}_{i=1}^{s}\sim p(D),
    \]
    where \(p(D)\) is the prior predictive distribution over
    whole datasets. A row permutation \(\pi\) and context size
    \(n\) define
    \[
    \begin{aligned}
        D_{\mathrm{ctx}}
        &=\{(x_{\pi(i)},y_{\pi(i)})\}_{i=1}^{n},\\
        D_{\mathrm{qry}}
        &=\{(x_{\pi(i)},y_{\pi(i)})\}_{i=n+1}^{s}.
    \end{aligned}
    \]
    There are \(s\) total rows and \(s-n\) query rows.
    The presentation's task loss is
    \[
        \mathcal{L}(\phi;D,\pi,n)
        =
        -\sum_{i=n+1}^{s}
        \log q_\phi
        \left(y_{\pi(i)}\mid x_{\pi(i)},D_{\mathrm{ctx}}\right).
    \]
    The overall objective averages this loss over datasets,
    permutations, and context sizes. This is a general PFN
    description; the local EPIT regression loss is specified in
    Section~\ref{sec:tabicl-regression}.

    \item \textbf{A concrete CorrPFN training task.}
    The reported run uses \(s=1{,}024\) synthetic rows.
    For an illustrative allowed split, \(n=700\) rows supply
    context and \(324\) supply queries; the actual context size
    varies across training tasks. For an informed task, the model
    receives all \(1{,}024\) rows of the standardized
    \(21\)-column input table and only the \(700\) context labels.
    It predicts distributions for the remaining \(324\) responses.
    Their generated labels enter CRPS, whose gradients update
    \(\phi\). New tasks repeat the procedure with different
    observations and sampled mechanisms.

    \item \textbf{Amortized inference.} The work of learning how to
    use different contexts is paid during pretraining and reused
    afterward. Amortized refers to spreading that learning
    cost across later prediction tasks. A new context changes
    internal representations and predictions without a gradient
    update to \(\phi\).

    \item \textbf{The corresponding real-data prediction.}
    At final evaluation, the \(608\) development rows and their
    EPIT labels supply context; the \(152\) final-test rows supply
    query features. The selected transformer applies the learned
    prediction procedure with fixed \(\phi\). Real query labels
    are used afterward for evaluation. The analytical rules have
    already influenced pretraining and are not added as a separate
    correction to these final predictions.

    \item \textbf{Why the approximation sign matters.} Finite
    training data, limited model capacity, imperfect optimization,
    and restrictions in the output distribution can all create
    approximation error. A mismatch between the synthetic prior
    and real corrosion data is a further issue. Approximating a
    synthetic-prior predictive distribution does not establish
    that the distribution is correct for the real dataset.

    \item \textbf{A statistical view without explicit posterior
    computation.} With \(\phi\) fixed, the PFN is a prediction rule
    whose output still depends on the observed sample \((x,y)\).
    Its behavior can therefore also be studied through how its
    predictions change across samples. Bayesian motivation does
    not imply automatic statistical consistency for arbitrary
    data distributions; see
    \href{https://proceedings.mlr.press/v202/nagler23a.html}
    {Nagler (2023), \textit{Statistical Foundations of Prior-Data
    Fitted Networks}}.
\end{notes}


\subsection{Where the Generic Synthetic Tasks Come From}
\label{sec:generic-synthetic-tasks}

So far, task sampling has been treated as an available operation.
This subsection explains how random functions and dependency
structures supply varied training problems. It describes the
generic generator first, then identifies how the informed branch
uses related components with the physical EPIT inputs.

\begin{notes}
    \item \textbf{The prior is a generator.} A configured prior
    produces many tasks with different sampled mechanisms,
    observations, and noise. It is not one fixed synthetic table.
    The generator's settings control the distribution from which
    the task mechanisms \(\theta\) are drawn.

    \item \textbf{MLP-generated tasks.} A randomly configured
    multilayer perceptron (MLP) creates nonlinear relationships.
    In the generic generator, one possible construction transforms
    latent root variables into intermediate variables, then selects
    features and a target from them:
    \[
        z\longrightarrow h\longrightarrow(X,y).
    \]
    Here, \(z\) and \(h\) belong to the synthetic generator, not
    the prediction transformer.

    \item \textbf{The two neural networks have different jobs.}
    A generator MLP creates a relationship from which observations
    are sampled. Its particular random weights help define
    \(\theta\). The TabICL transformer learns across those generated
    relationships by updating \(\phi\). Randomly changing the
    generator MLP creates another prediction problem; training
    the transformer improves its ability to solve such problems
    from their labeled examples.

    \item \textbf{Direct MLP or tree construction.} Another form
    generates features and then maps them to a target:
    \[
        X=z,\qquad y=f_\theta(X).
    \]
    Tree-based generators supply split-based or piecewise
    relationships. The mixed generator samples between the MLP
    and tree alternatives.

    \item \textbf{What SCM means.} A structural causal model
    describes variables using a directed graph without cycles
    and one equation for each variable. Following slide~15,
    \[
        v_j=f_j(\mathrm{pa}_j,\epsilon_j),
    \]
    where \(\mathrm{pa}_j\) denotes the parents of variable \(v_j\)
    in the graph and \(\epsilon_j\) denotes noise.
    Some generated variables can become observed features and
    another can become the target. This construction supplies
    synthetic dependencies; it does not establish causal effects
    in the real EPIT data.

    \item \textbf{Why latent variables are useful in the generic generator.}
    A shared hidden input can influence two observed columns and
    make them correlated. Other paths can make the response depend
    on interactions or leave some observed columns uninformative.
    Varying these constructions teaches the transformer to handle
    different statistical relationships. The graph notation
    explains this source of variation; it does not identify an
    unobserved physical cause in the corrosion dataset.

    \item \textbf{The informed target head uses a specific simpler path.}
    In the reported informed branch, physical alloy, environment,
    and method features are generated first. A random MLP or tree
    maps their standardized values directly to a response, with
    \texttt{is\_causal=False}. That response is combined with the
    analytical EPIT contribution. Thus, the generic latent-variable
    example above is background for the generator family; the
    informed implementation does not infer a causal graph from
    the real EPIT data or replace its physical inputs with latent
    features. Section~\ref{sec:informed-scm-target} gives the
    exact target construction.

    \item \textbf{Diversity and preprocessing.} Generic tasks can
    contain interactions, correlated or irrelevant inputs,
    different feature counts, and different context/query sizes.
    Depending on configuration, preprocessing can standardize,
    transform, permute, remove unusable columns, and pad inputs.
    Invalid tasks are rejected and sampled again.

    \item \textbf{Which probabilities apply to the reported model.}
    Trial~56 routes to generic tasks with probability \(0.25\),
    using MLP/tree probabilities of \(0.7/0.3\) within that branch.
    Its informed branch has routing probability \(0.75\) and
    uses MLP/tree probabilities of \(0.5/0.5\).
    These are sampling probabilities, not exact counts required
    in each batch. Some routing choices are shared by subgroups
    of tasks, as explained in Section~\ref{sec:generation-order}.

    \item \textbf{Why an informed prior is useful to investigate.}
    Generic columns have no built-in identity as chromium,
    chloride, or temperature. Their targets also have no fixed
    corrosion meaning. CorrPFN changes the task distribution
    so that part of the pretraining reflects the structure of the
    EPIT problem. The following sections define that construction.
\end{notes}

\subsection{How TabICL Uses the Context}
\label{sec:tabicl-context-processing}

We have described the prediction target, the training procedure,
and the source of synthetic examples. We now turn to the
computation that implements the learned function. Attention and
embeddings describe how the network processes its inputs; they
are not extra steps in the Bayesian reference calculation.

\begin{notes}
    \item \textbf{PFN, TabPFN, and TabICL.} PFN names the general
    approach. The presentation uses TabPFN to explain a tabular
    architecture. This project uses TabICL, which follows the
    in-context prediction idea with a different tabular architecture.
    The presentation's exact TabPFN layer counts and simple
    row-encoding example should therefore not be read as the
    configuration of CorrPFN.

    \item \textbf{The local TabICL computation.} The model first
    builds representations of cells using information from their
    columns. It then combines features within each row, and finally
    uses labeled context rows to predict query rows:
    \[
    \begin{gathered}
        \text{column representations}
        \longrightarrow \text{within-row feature interactions}\\
        \longrightarrow \text{in-context prediction}.
    \end{gathered}
    \]
    An \emph{embedding} or \emph{representation} is a learned
    vector used internally by the model.

    \item \textbf{Read the stages using one EPIT query.}
    Column processing represents values such as chromium amount
    using information from the corresponding column; the local
    target-aware components can also use context labels.
    Within-row processing combines the alloy composition with
    temperature, chloride, pH, and method for that observation.
    In-context processing then uses labeled examples to determine
    what response is plausible for the query. These are learned
    representations rather than explicit corrosion-rule terms.

    \item \textbf{Features and labels have different roles.}
    Context rows provide both their features and known targets.
    Query rows provide features only. Target information is
    supplied for the context; the unknown query labels are not
    input values. The model produces the query predictions from
    these representations.
    In the local in-context module, encoded context labels are
    added to the context-row representations before transformer
    processing. Unknown query labels receive no such encoding.

    \item \textbf{Rows form a set.} Context rows are examples,
    not positions in a time sequence. Changing their order should
    not change the prediction meaning. This is different from
    permuting feature columns: the EPIT columns have fixed roles
    and remain in their specified order.

    \item \textbf{Attention combines useful information.}
    In the basic attention example from slides~22--23, a receiving
    vector \(h^\ast\) is compared with source vectors \(h_j\).
    Learned projections produce
    \[
        q_\ast=W_Qh^\ast,\qquad
        k_j=W_Kh_j,\qquad
        v_j=W_Vh_j.
    \]
    The query vector \(q_\ast\) describes what information is
    being sought, the key \(k_j\) helps determine relevance, and
    the value \(v_j\) carries the information to combine.

    \item \textbf{Scores become weights.} In that basic example,
    \[
        s_j=\frac{q_\ast^\top k_j}{\sqrt{d}},
        \qquad
        \alpha_j=\frac{\exp(s_j)}{\sum_i\exp(s_i)},
        \qquad
        o_\ast=\sum_j\alpha_jv_j.
    \]
    Here, \(d\) is the vector dimension, the weights
    \(\alpha_j\) are nonnegative and sum to one, and \(o_\ast\)
    is the combined information. Different queries can place
    different weights on the available information. These
    equations explain basic attention; the local V2 components
    also use query-aware attention scaling.

    \item \textbf{An intuition for attention on corrosion data.}
    A query could benefit from examples with comparable compositions,
    similar environments, or informative material--environment
    contrasts. Attention provides a learned way to combine relevant
    internal information. Its similarity scores operate on learned
    vectors, so they need not equal raw composition distances or
    nearest-neighbour weights. This is an explanation of what the
    architecture can represent, not a measured interpretation of
    particular attention heads in the reported model.

    \item \textbf{Attention weights are not posterior mechanism
    probabilities.} The \(\alpha_j\) values weight internal
    vectors. They are not explicit values of
    \(\pr(\theta\mid x,y)\). Likewise, the attention vector
    \(q_\ast\) is different from the predictive distribution
    \(q_\phi\).

    \item \textbf{What remains learned after pretraining.}
    The projection weights and other transformer parameters are
    part of \(\phi\). New context inputs change the vectors,
    attention scores, and output, while these parameters stay
    fixed.

    \item \textbf{Why fixed weights still allow adaptation.}
    A fixed prediction procedure can produce different answers
    when its inputs change. For comparison, the least-squares
    formula is fixed, but different \((X,y)\) produce different
    fitted coefficients. TabICL learns its own computation across
    tasks: changing the context changes its internal vectors and
    predicted distribution without a gradient update to \(\phi\).
    This analogy explains dependence on the data; it does not mean
    that TabICL necessarily computes least-squares coefficients.
\end{notes}

\subsection{Regression, Loss Functions, and Predictive Uncertainty}
\label{sec:tabicl-regression}

Context processing produces a representation for each query.
We now specify what prediction is read from it and how that
prediction is scored during training. The Gaussian and Laplace
examples below explain general connections between distributions
and familiar losses. CorrPFN specifically uses the quantile
output and CRPS described afterward.

\begin{notes}
    \item \textbf{Why the loss matters.} The training loss specifies
    what a good prediction means. A point-prediction loss assesses
    a selected value. A distributional loss assesses the
    predictive distribution.

    \item \textbf{Squared error and a Gaussian likelihood.}
    Slide~11 considers
    \(q_\phi(\cdot\mid x^\ast,x,y)=
    \operatorname{Normal}(\mu_\phi,\sigma^2)\).
    When \(\sigma^2\) is fixed,
    \[
        -\log q_\phi(y^\ast\mid x^\ast,x,y)
        =
        \frac{(y^\ast-\mu_\phi)^2}{2\sigma^2}
        +\text{constant}.
    \]
    Minimizing this negative log-likelihood is therefore
    equivalent to minimizing mean squared error. The predictive
    mean \(\mu_\phi\) depends on the context and query.

    \item \textbf{Absolute error and a Laplace likelihood.}
    With a fixed scale \(b>0\), the same slide gives
    \[
    \begin{aligned}
        q_\phi(\cdot\mid x^\ast,x,y)
        &=\operatorname{Laplace}(\mu_\phi,b),\\
        -\log q_\phi(y^\ast\mid x^\ast,x,y)
        &=\frac{|y^\ast-\mu_\phi|}{b}+\text{constant}.
    \end{aligned}
    \]
    This connects mean absolute error to a different assumed
    predictive noise distribution.

    \item \textbf{Learning the spread as well as the centre.}
    Slide~12 lets a Gaussian PFN predict both \(\mu_\phi\) and
    \(\sigma_\phi^2\), with loss
    \[
        \mathcal{L}(\phi)
        =
        \mathbb{E}\left[
            \frac{1}{2}\log\sigma_\phi^2
            +\frac{(y^\ast-\mu_\phi)^2}{2\sigma_\phi^2}
        \right]+\text{constant}.
    \]
    The model must balance accurate predictions with a suitable
    spread. The logarithmic term prevents it from reducing the
    error penalty simply by making the variance arbitrarily large.
    This is an explanatory example from the presentation.

    \item \textbf{The actual EPIT model uses quantiles.}
    CorrPFN uses a different regression head: it predicts
    \(999\) quantiles at probability levels
    \(0.001,0.002,\ldots,0.999\).
    A quantile gives a response value associated with a cumulative
    probability; the \(0.5\) quantile is the median.
    These outputs are used to construct a predictive distribution
    without restricting it to the Gaussian example above.

    \item \textbf{Reading a quantile prediction.}
    Suppose, illustratively, that after conversion to millivolts
    the \(0.10\), \(0.50\), and \(0.90\) quantiles are
    \(100\), \(250\), and \(500\)~mV. The first means the model
    assigns approximately \(10\%\) probability to EPIT at or below
    \(100\)~mV; the middle value is its median prediction.
    The \(999\) outputs describe one predictive distribution for
    a query. They are not \(999\) sampled mechanisms or separate
    models. These example values are not experimental results.

    \item \textbf{The actual training loss is CRPS.}
    The local regression path minimizes the
    \emph{continuous ranked probability score}, averaged over
    synthetic query targets. To express it, write
    \(F_\phi(t\mid x^\ast,x,y)\) for the cumulative distribution
    associated with \(q_\phi\): its predicted probability that
    the response is at most \(t\). In the loss convention,
    \begin{equation}
        \operatorname{CRPS}(F_\phi,y^\ast)
        =
        \int_{-\infty}^{\infty}
        \left[
            F_\phi(t\mid x^\ast,x,y)
            -\mathbf{1}\{y^\ast\leq t\}
        \right]^2\,dt.
        \label{eq:crps-loss}
    \end{equation}
    The indicator is one when the observed outcome is at most
    \(t\), and zero otherwise. Lower CRPS is better.

    \item \textbf{An intuitive threshold calculation.}
    On an illustrative millivolt scale, suppose the observed EPIT
    is \(250\)~mV. At threshold \(t=300\)~mV, the event
    \emph{EPIT is at most \(300\)~mV} occurred, so the indicator is
    one. Predicted probabilities of \(0.8\) and \(0.2\) give
    squared errors \((0.8-1)^2=0.04\) and
    \((0.2-1)^2=0.64\), respectively. CRPS integrates these
    comparisons across thresholds; it balances concentration
    with assigning probability to plausible outcomes. Actual
    synthetic training uses the generated standardized target scale.

    \item \textbf{Statistical purpose of CRPS.} It compares
    predicted cumulative probabilities with the observed
    outcome across thresholds. Like log loss, ideal expected
    CRPS favors the true conditional distribution, under its
    usual finite-moment assumptions. Thus, changing from the
    slides' log loss to CRPS still supports learning a
    predictive distribution. See
    \href{https://doi.org/10.1198/016214506000001437}
    {Gneiting and Raftery (2007)}.

    \item \textbf{How this completes the PFN training argument.}
    At each threshold, expected squared error is minimized by
    predicting the true conditional probability of the event.
    Integrating over thresholds therefore favors the true
    conditional CDF, under the stated moment assumptions.
    This is the connection to Section~\ref{sec:pfn-pretraining}:
    CRPS, like log loss, lets generated query outcomes teach a
    conditional distribution without exact posterior calculations
    as training labels. The target distribution is the one induced
    by the configured synthetic tasks; its accuracy for real
    corrosion remains an empirical question.

    \item \textbf{What is reported as the prediction.} The
    reported EPIT point estimate is the median of the predicted
    distribution. Target scaling is fitted on context labels,
    and predictions are transformed back to millivolts.
    A distributional training objective does not, by itself,
    establish that prediction intervals are calibrated on real
    corrosion data.

    \item \textbf{How a standardized prediction returns to millivolts.}
    The synthetic response has no literal EPIT unit. For real
    inference, the estimator fits target standardization on the
    context EPIT labels and reverses it for the output. If the
    context mean were \(200\)~mV, its standard deviation
    \(300\)~mV, and the predicted standardized median \(0.5\),
    the returned value would be \(200+300(0.5)=350\)~mV.
    These numbers illustrate scaling, not a reported prediction.
    Fitting this context scaling does not update transformer weights.

    \item \textbf{Keep training, selection, and reporting distinct.}
    Synthetic query CRPS supplies the gradients that train
    \(\phi\). Mean development-fold Spearman selects the prior
    configuration and checkpoint. Final Spearman, MAE, RMSE,
    and \(R^2\) assess the reported point predictions.
    These quantities serve different stages of the workflow;
    Spearman is not the transformer's gradient-training loss.
    Likewise, a fixed transformer can output a spread of plausible
    responses without sampling a posterior over its own weights.
\end{notes}

\subsection{Which Implementation and Training Scope This Project Uses}
\label{sec:tabicl-provenance}

\begin{notes}
    \item \textbf{Local architecture.} CorrPFN and the local
    generic baseline use the same configured transformer
    architecture, including V2 components such as target-aware
    column embeddings and query-aware scalable softmax.
    Their saved architecture configurations and parameter
    shapes match.

    \item \textbf{Generic generator and training framework.}
    The generic MLP/tree generators and Stage~1 training
    framework derive from TabICLv1. The project adds the
    informed prior and the continuous-target training path,
    using the V2 quantile/distribution implementation.

    \item \textbf{Pretrained weights are a separate choice.}
    Both local models are trained from scratch.
    The authors' released pretrained TabICL V2 weights are
    used only for the external reference model. Using V2
    architectural components locally does not mean reproducing
    the authors' full V2 pretraining recipe.

    \item \textbf{Version labels refer to different things.}
    A project label such as \texttt{v7} identifies a revision of
    the CorrPFN experiment pipeline. TabICL V1 and V2 identify
    upstream model generations. A software-package version is
    a third, separate label.

    \item \textbf{Stage~1 scope.} The reported local runs use
    synthetic tables of \(1{,}024\) rows and only Stage~1 of the
    original curriculum. Later stages extend long-context
    capability. Earlier attempts to extend training did not
    improve loss in the implementations tested at that time;
    this does not rule out improvements from newer versions.

    \item \textbf{Dataset adaptation and inference are separate.}
    CorrPFN's prior is adapted to the EPIT task. Real development
    labels influence rule calibration and model selection,
    while transformer weight optimization uses synthetic
    tasks. The absence of gradient updates during final
    inference does not mean that the whole project used no
    real labels during development.

    \item \textbf{Implementation references.}
    The architecture and loss descriptions above can be traced
    to \path{src/tabicl/model/tabicl.py},
    \path{src/tabicl/model/learning.py},
    \path{src/tabicl/model/quantile_dist.py}, and
    \path{src/tabicl/train/run.py}.
    The upstream model references are
    \href{https://arxiv.org/abs/2502.05564}{Qu et al.,
    \textit{TabICL} (2025)} and
    \href{https://arxiv.org/abs/2602.11139}{Qu et al.,
    \textit{TabICLv2} (2026)}.
\end{notes}

\subsection{Figures for the Presentation, Described in Words}
\label{sec:presentation-figure-description}

\begin{notes}
    \item \textbf{Main figure: where the labeled data enter.}
    Draw two horizontal lanes with the headings
    \emph{Fit a model for this task} and
    \emph{Use a pretrained PFN for this task}.
    Use the same small example table and query row in both
    lanes so that the difference is the computation.

    \item \textbf{Upper lane: conventional fitting.}
    Show the labeled table \((x,y)\) entering a box marked
    \emph{Fit model parameters}. An orange update loop produces
    \(\widehat\theta\). Place a separate prediction box to its
    right: it receives the fitted model and \(x^\ast\), then
    outputs a prediction. Add a small note that another task
    normally requires another fit or adaptation.

    \item \textbf{Lower lane, first part: PFN pretraining.}
    Show several small synthetic tables generated from different
    \(\theta\sim\pr(\theta)\). Arrows lead to a transformer
    training box that updates \(\phi\). Show a sampled query
    target \(y^\ast\) going to the loss calculation, not to the
    transformer input. Put this part in a shaded area labeled
    \emph{Pretraining across many tasks}.

    \item \textbf{Lower lane, second part: PFN inference.}
    After a clear divider, show the real labeled context
    \((x,y)\) and query \(x^\ast\) entering the frozen transformer
    together. Put a lock symbol on \(\phi\). Its output is
    \(q_\phi(y^\ast\mid x^\ast,x,y)\), drawn as a small
    predictive distribution. There is no parameter-update loop
    in this part.

    \item \textbf{Use consistent visual labels.} Orange marks
    parameter updates; a lock marks fixed parameters.
    Use the same colour for the real context in both lanes.
    Mark the unknown query target with a question mark.
    Label the main distinction:
    \emph{New labels are used for parameter fitting above,
    and as context inputs below}.

    \item \textbf{What the main figure should communicate.}
    Both approaches use labeled data and can produce
    probabilistic predictions. The PFN has already learned how
    to use different contexts during pretraining.
    The figure should not imply that PFNs require no training,
    or that conventional models must update their weights for
    every individual query.

    \item \textbf{Optional second figure: MAP versus Bayesian
    prediction versus PFN.} Use a one-dimensional synthetic
    example, as in slide~9: input on the horizontal axis and
    response on the vertical axis. Show the same black context
    points and a blue query position in three panels.
    The MAP panel shows one selected mechanism's curve.
    The Bayesian panel shows several plausible curves with
    thickness indicating posterior weight, together with their
    combined predictive distribution at the query.
    The PFN panel shows a learned approximation to that
    distribution, labeled \(q_\phi\approx\pr\).
    Mark the whole example as illustrative.

    \item \textbf{A useful animation.} Add one new labeled
    context point. In the Bayesian panel, the relative weights
    of mechanisms change. In the PFN panel, the prediction
    changes after another forward computation while the lock
    on \(\phi\) stays visible. This makes context adaptation
    without weight updates concrete.
\end{notes}

\section{How the Synthetic Training Tasks Are Generated}
\label{sec:informed-prior}

\subsection{Starting Point and Order of Operations}
\label{sec:generation-order}

\begin{notes}
    \item \textbf{What is generated.} One synthetic task contains
    input rows and their generated responses. Following the
    presentation's notation, \(s\) is the total number of rows and
    \(n\) is the number later used as labeled context. For the
    informed branch, the raw input table is
    \[
        X_{\mathrm{raw}}\in\mathbb{R}^{s\times21}.
    \]
    Its columns follow the dataset schema in
    Section~\ref{sec:dataset-features}.

    \item \textbf{What must already be available.} Before transformer
    pretraining starts, the generator needs the saved feature
    profiles, calibrated analytical-rule coefficients, and
    rule-sampling probabilities. The coefficient-fitting procedure
    is covered in Section~4. The steps below describe how these
    prepared inputs are used to generate training tasks.

    \item \textbf{Choose the task branch first.} The hybrid prior
    uses an informed task with probability \(\rho\), and a generic
    task with probability \(1-\rho\). Using the presentation's
    \(p(D)\) notation for a distribution over datasets,
    \begin{equation}
        p(D)=\rho\,p_{\mathrm{informed}}(D)
             +(1-\rho)\,p_{\mathrm{generic}}(D).
        \label{eq:hybrid-task-distribution}
    \end{equation}
    The setting \(\rho\) is
    \texttt{informed\_prior\_ratio}. It controls how often the
    informed branch is used.

    \item \textbf{Order within an informed task.} Sample alloy
    compositions; sample environmental conditions and test methods;
    keep raw and standardized copies of the resulting table;
    generate an SCM response and an analytical EPIT response;
    combine the responses; then supply context and query rows
    for transformer training.

    \item \textbf{The generic branch.} Generic tasks use the
    construction described in Section~\ref{sec:generic-synthetic-tasks}.
    Their MLP/tree probabilities remain \(0.7/0.3\).
    They use their existing preprocessing and receive neither the
    empirical EPIT feature table nor an analytical EPIT contribution.

    \item \textbf{How this is batched in the implementation.}
    Routing and several generator settings are sampled for a
    subgroup of datasets. Datasets in that subgroup share those
    choices, while their observations and noise are sampled
    separately. Thus, the routing choices are not independent
    for every dataset in a batch.
\end{notes}

\subsection{Saved Feature Profiles and Their Data Sources}
\label{sec:feature-profile-provenance}

\begin{notes}
    \item \textbf{Composition profile.}
    \texttt{epit\_dataset\_v1} stores \(403\) distinct reported
    composition templates. Each template contains the complete
    \(24\)-element source composition and its alloy-family label.
    The final informed branch draws from the \(298\) Fe-alloy
    templates and \(17\) Ni--Cr--Mo templates. Several related
    templates can belong to one composition group used for splitting.

    \item \textbf{Environmental profile.}
    \texttt{epit\_fe\_nicrmo\_features\_v1} stores the
    \(565\) Fe/Ni--Cr--Mo experimental records used to sample
    temperature, chloride concentration, pH, and test method.
    It also records the numeric imputation means and method encoding.

    \item \textbf{Scope of the two alloy families.}
    The Fe/Ni--Cr--Mo restriction applies to these informed
    feature draws and to analytical-rule calibration. The
    transformer evaluation still includes all alloy classes
    in the dataset.

    \item \textbf{Where the profile information came from.}
    The profiles were built using feature information from the
    full \(760\)-row usable corpus. The imputation means and
    complete method mapping also use that corpus. They were
    not fitted only on the \(608\) development rows.

    \item \textbf{What this means for the final-test boundary.}
    The profiles contain no EPIT values, although target
    availability was used to select the usable rows. They do
    expose the generator to final-test feature information,
    including compositions assigned to final-test groups.
    The composition split therefore separates labeled examples;
    it does not establish that final-test compositions were
    unseen during synthetic feature generation.

    \item \textbf{Use during training.} These profiles are fixed,
    versioned assets. The generator loads them directly; it does
    not reopen the source Excel workbook for each task.
\end{notes}

\subsection{Sampling and Perturbing an Alloy Composition}
\label{sec:composition-generator}

\begin{notes}
    \item \textbf{Choose a family for each row.} Sampling
    probabilities follow the numbers of distinct templates:
    \[
        \pr(\mathrm{Fe})=\frac{298}{315}\approx0.9460,
        \qquad
        \pr(\mathrm{Ni\!-\!Cr\!-\!Mo})
            =\frac{17}{315}\approx0.0540.
    \]
    A template is then sampled uniformly within the chosen
    family. These frequencies count distinct compositions,
    rather than repeated experimental measurements.

    \item \textbf{Perturb the reported element amounts.}
    Let \(c_{ij}\) be the source amount of element \(j\) in
    the template selected for row \(i\). For positive amounts,
    \begin{equation}
        \widetilde c_{ij}=c_{ij}\exp(\epsilon_{ij}),
        \qquad
        \epsilon_{ij}\sim\operatorname{Normal}(0,\tau^2).
        \label{eq:composition-perturbation}
    \end{equation}
    The perturbations are sampled independently before
    composition closure. Larger \(\tau\) produces stronger
    changes around a template.
    Its setting is \texttt{pitting\_composition\_perturb\_strength}.

    \item \textbf{Preserve absent elements.} Zero and unreported
    source amounts remain zero. The perturbation changes the
    amounts of elements already present in the template.

    \item \textbf{Close the full composition to \(100\) wt.\%.}
    All \(24\) source elements take part in the normalization:
    \begin{equation}
        c_{ij}^{\mathrm{closed}}
        =
        100\frac{\widetilde c_{ij}}
                 {\sum_{e=1}^{24}\widetilde c_{ie}}.
        \label{eq:composition-closure}
    \end{equation}
    Only afterward are the \(17\) model-visible elements selected.
    Their sum can be below \(100\%\), because the omitted elements
    can account for part of the alloy.

    \item \textbf{Meaning of \(\tau=0\).} With no perturbation,
    the generator reuses a template after full-composition
    closure. Its values can still differ slightly from the
    original reported values if their total was not exactly
    \(100\) wt.\%.

    \item \textbf{Composition structure retained by this procedure.}
    The templates retain the dominant elements, typical amounts,
    and patterns of elements occurring together. Closure keeps
    the full vector consistent with a total of \(100\) wt.\%.
    The final branch has no extra shared block coupling or
    Dirichlet redraw; those belong to other generator variants
    discussed in earlier development.
\end{notes}

\subsection{Sampling the Environment and Test Method}
\label{sec:environment-generator}

\begin{notes}
    \item \textbf{Sample one complete experimental tuple.}
    Independently of the composition, choose one of the
    \(565\) environmental records uniformly with replacement.
    Copy its temperature, chloride concentration, pH, and
    test method together:
    \[
        (T_i,C_{\mathrm{Cl},i},\mathrm{pH}_i,a_i).
    \]
    Here, \(a_i\) is the encoded test-method category.

    \item \textbf{Which relationships are preserved.}
    Keeping the tuple intact preserves observed combinations
    of environmental conditions and methods. Independent
    composition sampling allows those conditions to be paired
    with a different alloy from the template bank.
    The resulting complete row need not have occurred in the
    real dataset.

    \item \textbf{Missing values and category codes.}
    Missing numeric environmental values use the means saved
    in the profile. Methods use the fixed lexicographically
    ordered mapping over \(52\) categories, including the
    missing-value category. This is the synthetic profile's
    mapping; context-fitted evaluation encodings are described
    in Section~\ref{sec:real-data-representation}.

    \item \textbf{Assemble the physical row.}
    Append the three environmental values and method code to
    the \(17\) element amounts. Every informed task uses this
    complete \(21\)-column profile, with the same feature order.
\end{notes}

\subsection{Keeping Raw and Standardized Feature Tables}
\label{sec:feature-standardization}

\begin{notes}
    \item \textbf{Keep the physical values.}
    \(X_{\mathrm{raw}}\) retains compositions in wt.\%,
    temperature in degrees Celsius, chloride concentration in
    molar units, pH, and the method identifiers.

    \item \textbf{Make a standardized copy.}
    Let \(\mu_j\) and \(\sigma_j\) be the population mean and
    standard deviation of column \(j\) across all \(s\) rows
    in the generated task. Then
    \begin{equation}
        X_{\mathrm{std},ij}
        =
        \begin{cases}
            \displaystyle
            \frac{X_{\mathrm{raw},ij}-\mu_j}{\sigma_j},
                & \sigma_j>10^{-6},\\[6pt]
            0,  & \sigma_j\leq10^{-6}.
        \end{cases}
        \label{eq:feature-standardization}
    \end{equation}

    \item \textbf{Use each copy where it is needed.}
    \begin{equation}
        X_{\mathrm{raw}}\longrightarrow
        \begin{cases}
            g_k(X_{\mathrm{raw}}),
                & \text{analytical EPIT response},\\
            f_\theta(X_{\mathrm{std}}),
                & \text{SCM response},\\
            X_{\mathrm{std}},
                & \text{transformer inputs}.
        \end{cases}
        \label{eq:informed-pipeline}
    \end{equation}
    The analytical rule needs raw values for its thresholds,
    concentration terms, and method lookup. Standardized
    features give the random SCM comparable input scales.
    The SCM and transformer receive the same standardized
    physical columns.

    \item \textbf{Preserve the schema.} Constant columns become
    zero and keep their positions. Feature permutation is
    disabled. The model receives numeric columns without
    textual feature names, so stable positions preserve their
    roles. The informed branch bypasses the general
    \texttt{Reg2Cls} transformation.

    \item \textbf{Standardization of response signals.}
    In the equations below, \(\mathcal{Z}\) acts on one
    row-wise signal \(v\) across the generated task:
    \[
        \mathcal{Z}(v)_i
        =
        \frac{v_i-\overline v}
             {\max(\operatorname{sd}_{\mathrm{pop}}(v),10^{-6})}.
    \]
    Exactly constant signals become zero. The denominator floor
    handles nearly constant signals. These statistics are
    recomputed for each synthetic task; they are distinct
    from the context-fitted states used for real-data evaluation.
\end{notes}

\subsection{Generating the SCM Response}
\label{sec:informed-scm-target}

\begin{notes}
    \item \textbf{Choose the stochastic target head.}
    Within the informed branch,
    \[
        \pr(\mathrm{MLP}\mid\mathrm{informed})=p_{\mathrm{MLP}},
        \qquad
        \pr(\mathrm{tree}\mid\mathrm{informed})=1-p_{\mathrm{MLP}}.
    \]
    The search setting \texttt{mlp\_prob} controls this choice.
    It is separate from \(\rho\), which selects the task branch.

    \item \textbf{Generate one response per row.}
    The selected head maps the standardized physical inputs to
    a response:
    \begin{equation}
        y_{\mathrm{SCM}}
        =
        f_\theta^{\mathrm{MLP/tree}}(X_{\mathrm{std}}).
        \label{eq:scm-target}
    \end{equation}
    This is a direct feature-to-target construction, configured
    as \texttt{is\_causal=False}. The input table stays fixed.
    Its statistical relationships can be nonlinear, but it is
    not an identified causal model of corrosion.

    \item \textbf{MLP head.} Its depth, width, activation,
    initialization, dropout structure, and noise are sampled
    through the existing MLP-SCM generator. The final layer
    returns one target value for each row.

    \item \textbf{Tree head.} Each layer fits a regression tree
    or ensemble to random Gaussian targets using its current
    input table, then predicts on that table. These outputs
    become the next layer's inputs, with additive noise in
    subsequent layers. This produces random piecewise
    relationships. The fitted targets in this construction are
    synthetic random values; they are not real EPIT measurements.
\end{notes}


\subsection{Selecting the Analytical EPIT Rule}
\label{sec:analytical-target-selection}

\begin{notes}
    \item \textbf{Select one of seven rule families.} Let \(s_k\)
    be the saved development-validation score of rule \(k\).
    All seven reported scores are positive, and the sampler uses
    \[
        p_k=\frac{s_k}{\sum_{\ell=1}^{7}s_\ell}.
    \]
    Better-scoring rules are sampled more often. These are
    sampling probabilities chosen from validation scores, not
    Bayesian posterior probabilities. Section~4 explains how
    the scores are obtained.

    \item \textbf{Use one rule across the task.}
    The selected rule \(g_k\) acts on the raw physical table:
    \begin{equation}
        y_{\mathrm{EPIT}}=g_k(X_{\mathrm{raw}}).
        \label{eq:epit-target}
    \end{equation}
    Every row uses the same rule family and its saved
    coefficient vector. A rule family describes a mathematical
    response formula; it is distinct from an alloy family.

    \item \textbf{What is fixed and what changes.}
    The numerical constants inside the formulas and the saved
    coefficient vectors stay fixed across training tasks.
    The selected rule, feature values, and term-standardization
    statistics can change. The method-aware rule also draws
    fresh method offsets.

    \item \textbf{Additional generator information.}
    The Fe/Ni threshold rule uses the sampled template's
    alloy-family label. This label is internal metadata and
    is not an extra transformer input. The compact notation
    \(g_k(X_{\mathrm{raw}})\) leaves this metadata and the
    random method offsets implicit.

    \item \textbf{Role of the formulas.}
    The rules add selected corrosion-related relationships to
    the synthetic response. They are heuristic scores with
    calibrated relative weights, not equations that directly
    return physical EPIT values in millivolts. Their complete
    terms are given below.
\end{notes}

\subsection{Material Terms in the Analytical Rules}
\label{sec:analytical-material-terms}

\begin{notes}
    \item \textbf{Notation and input units.}
    Element symbols denote nonnegative raw weight-percent
    amounts. \(\sigma\) denotes the logistic sigmoid,
    \(\log\) is the natural logarithm, and \(\mathcal{Z}\)
    is the task-wise standardization defined in
    Section~\ref{sec:feature-standardization}.
    The numerical constants below assume these input units.

    \item \textbf{Combine the Mo and W contributions.}
    \[
        Q=\mathrm{Mo}+0.55\mathrm{W}.
    \]
    This quantity is used inside the material scores; the
    original Mo and W columns remain separate model inputs.

    \item \textbf{Linear PREN-inspired score.}
    \[
        R_{\mathrm{linear}}=\mathrm{Cr}+3.25Q.
    \]
    This gives positive material-resistance contributions to
    Cr, Mo, and W. It is not a complete conventional PREN
    formula: nitrogen is absent from the retained schema.

    \item \textbf{Cr--Mo/W synergy score.}
    \[
        R_{\mathrm{synergy}}
        =\mathrm{Cr}+3.25Q+\sqrt{\mathrm{Cr}\,Q}.
    \]
    The extra term represents a joint contribution from
    chromium and Mo/W.

    \item \textbf{Common chromium-threshold score.}
    \[
        R_{\mathrm{threshold}}
        =\sigma\!\left(\frac{\mathrm{Cr}-12}{2}\right)
         +\log(1+Q).
    \]
    This uses a chromium threshold near \(12\) wt.\% and a
    Mo/W contribution with diminishing increments.

    \item \textbf{Family-dependent chromium-threshold score.}
    \[
    \begin{aligned}
        R_{\mathrm{Fe/Ni}}
        &=\sigma\!\left(\frac{\mathrm{Cr}-t_i}{2}\right)
          +\log(1+Q),\\
        t_i&=
        \begin{cases}
            12,&\text{Fe-alloy row},\\
            15,&\text{Ni--Cr--Mo row}.
        \end{cases}
    \end{aligned}
    \]
    The template metadata determine which threshold is used.

    \item \textbf{Turn the selected score into rule terms.}
    For the material score \(R_k\) used by rule \(k\),
    \begin{equation}
    \begin{aligned}
        m_k&=\mathcal{Z}(R_k),\\
        P_k&=\mathcal{Z}(\tanh m_k),\\
        I_k&=\mathcal{Z}\!\left(\sigma(-m_k)C\right).
    \end{aligned}
    \label{eq:rule-material-signals}
    \end{equation}
    \(P_k\) is the material-protection term.
    \(I_k\) combines material susceptibility with the chloride
    drive \(C\), defined next. It enters the relevant rules
    with a negative sign.
\end{notes}

\subsection{Environmental and Test-Method Terms}
\label{sec:analytical-environment-terms}

\begin{notes}
    \item \textbf{Chloride terms.}
    With \(C_{\mathrm{Cl}}\) in molar units,
    \[
        L=\mathcal{Z}\!\left(
            \log_{10}\max(C_{\mathrm{Cl}},10^{-12})
        \right),
        \qquad C=\sigma(L).
    \]
    The concentration floor allows zero reported chloride
    values to pass through the logarithm.

    \item \textbf{Combined environmental term.}
    \begin{equation}
        E=\sigma\!\left(
            \mathcal{Z}\!\left[
                0.075\mathcal{Z}(T)+0.925L
                +0.115\mathcal{Z}(|\mathrm{pH}-7.25|)
            \right]
        \right).
        \label{eq:rule-combined-environment}
    \end{equation}
    This combines temperature, logarithmic chloride, and
    departure from a reference pH of \(7.25\). The fixed
    internal weights give chloride the largest coefficient.

    \item \textbf{Separate temperature and acidic-pH terms.}
    \[
        H=\sigma\!\left(\frac{T-50}{10}\right),
        \qquad A=\sigma(6.5-\mathrm{pH}).
    \]
    \(H\) represents higher-temperature conditions, with its
    transition centered at \(50\,^{\circ}\mathrm{C}\).
    \(A\) represents an acidic-pH contribution.

    \item \textbf{Temperature--chloride interactions.}
    \[
        J_0=\sigma(\mathcal{Z}(T))C,
        \qquad
        J_1=HC.
    \]
    \(J_0\) uses temperature relative to the generated task;
    \(J_1\) uses the fixed high-temperature term.
    Both combine temperature with chloride.

    \item \textbf{Coupled material/environment breakdown.}
    \begin{equation}
        D_k=\sigma\!\left[
            0.65L+0.25\mathcal{Z}(H)
            +0.10\mathcal{Z}(J_1)-0.75m_k
        \right].
        \label{eq:rule-coupled-breakdown}
    \end{equation}
    Environmental contributions increase this breakdown
    score, while the material score opposes them.
    \(D_k\) is a rule term; \(D\) without a subscript denotes
    a dataset elsewhere in the report.

    \item \textbf{Random test-method term.}
    For a method-aware task, draw an offset for each encoded
    category:
    \[
        o_j\sim\operatorname{Normal}(0,1),
        \qquad j=0,\ldots,51.
    \]
    A row with category \(a_i\) uses \(v_i=o_{a_i}\), giving
    \begin{equation}
        M=\mathcal{Z}\!\left[\tanh\!\left(\mathcal{Z}(v)\right)\right].
        \label{eq:rule-method-term}
    \end{equation}
    Rows with the same method receive the same offset within
    a task. Both standardizations act across rows.
    These offsets are resampled for each task; they do not
    copy fitted real-data method effects.

    \item \textbf{How to interpret the environmental assumptions.}
    The rules offer different candidate relationships.
    In particular, the pH-departure term in \(E\) and the
    acidic term \(A\) express different assumptions.
    Their presence does not establish a universal pH law
    for every alloy and experimental range.
\end{notes}

\subsection{The Seven Complete Analytical Rules}
\label{sec:analytical-rule-families}

\begin{notes}
    \item \textbf{Common weighted form.}
    Each rule combines its listed terms:
    \begin{equation}
        g_k(X_{\mathrm{raw}})
        =b_k(X_{\mathrm{raw}})c_k,
        \qquad
        c_{k,j}\geq0,\qquad
        \sum_j c_{k,j}=1.
        \label{eq:analytical-rule-vector}
    \end{equation}
    The columns of \(b_k\) contain the signed terms in the
    order shown below. The entries of \(c_k\) are their
    saved calibrated weights. Nonnegative weights preserve
    the signs assigned to those terms; a zero coefficient
    makes its term inactive.

    \item \textbf{Linear PREN}
    (\texttt{pren\_linear}).
    Use \(R_{\mathrm{linear}}\) and
    \[
        b_k=(P_k,-\mathcal{Z}(E),-I_k).
    \]
    Material protection competes with the combined environment
    and susceptibility--chloride contributions.

    \item \textbf{Cr--Mo/W synergy}
    (\texttt{cr\_mow\_synergy}).
    Use \(R_{\mathrm{synergy}}\) and
    \[
        b_k=(P_k,-\mathcal{Z}(E),-I_k,-\mathcal{Z}(J_0)).
    \]
    This adds material synergy and a separate
    temperature--chloride contribution.

    \item \textbf{Threshold and saturation}
    (\texttt{threshold\_saturation}).
    Use \(R_{\mathrm{threshold}}\) and
    \[
        b_k=(P_k,-\mathcal{Z}(E),-I_k,-\mathcal{Z}(J_0)).
    \]
    Its environmental terms match the synergy rule, while
    its material score uses the chromium threshold and
    logarithmic Mo/W contribution.

    \item \textbf{Improved environment}
    (\texttt{improved\_environment}).
    Use \(R_{\mathrm{linear}}\) and
    \[
        b_k=(P_k,-\mathcal{Z}(C),-\mathcal{Z}(H),
                   -\mathcal{Z}(J_1),-\mathcal{Z}(A)).
    \]
    Chloride, high temperature, their interaction, and acidic
    pH have separate weights. The name identifies this
    candidate formula; it does not by itself imply better
    prediction performance.

    \item \textbf{Coupled breakdown}
    (\texttt{coupled\_breakdown}).
    Use \(R_{\mathrm{linear}}\) and
    \[
        b_k=(P_k,-\mathcal{Z}(D_k),-\mathcal{Z}(A)).
    \]
    The breakdown term combines environmental aggressiveness
    and material resistance before entering the weighted rule.

    \item \textbf{Fe/Ni chromium threshold}
    (\texttt{fe\_ni\_cr\_threshold}).
    Use \(R_{\mathrm{Fe/Ni}}\) and
    \[
        b_k=(P_k,-\mathcal{Z}(C),-\mathcal{Z}(H),
                   -\mathcal{Z}(J_1),-\mathcal{Z}(A)).
    \]
    This pairs the family-dependent material threshold with
    the separate environmental terms.

    \item \textbf{Method-aware}
    (\texttt{method\_aware}).
    Use \(R_{\mathrm{synergy}}\) and
    \[
        b_k=(P_k,-\mathcal{Z}(E),-I_k,-\mathcal{Z}(J_0),M).
    \]
    The method term adds task-specific variation associated
    with the measurement category. Its contribution can be
    positive or negative because \(M\) is centered.
\end{notes}

\subsection{Combining the SCM and Analytical Responses}
\label{sec:target-mixture}

\begin{notes}
    \item \textbf{Standardize, combine, and standardize again.}
    The final informed response is
    \begin{equation}
        y_{\mathrm{informed}}
        =
        \mathcal{Z}\!\left[
            (1-\lambda)\mathcal{Z}(y_{\mathrm{SCM}})
            +\lambda\mathcal{Z}(y_{\mathrm{EPIT}})
        \right],
        \qquad 0\leq\lambda\leq1.
        \label{eq:scm-epit-target}
    \end{equation}
    The setting \(\lambda\) is
    \texttt{informed\_target\_mix\_weight}.

    \item \textbf{Why standardize the components separately.}
    The SCM and analytical rule can produce different
    numerical scales. Separate standardization lets
    \(\lambda\) control their relative contribution without
    those arbitrary scales dominating the mixture.

    \item \textbf{Meaning of the endpoints.}
    For usable, nonconstant components, \(\lambda=0\) gives
    the SCM response, and \(\lambda=1\) gives the analytical
    response. Intermediate values combine stochastic
    relationships with the calibrated EPIT contribution.
    If the analytical response has a population standard deviation
    at or below \(10^{-6}\), its contribution is set to zero.

    \item \textbf{Different weights control different choices.}
    \(\rho\) selects the informed or generic task branch.
    Within an informed task, \(\lambda\) combines the two
    responses. Within the analytical response, \(c_k\)
    weights its individual terms. The rule probability \(p_k\)
    selects which analytical formula supplies that response.

    \item \textbf{The mixture is not a variance decomposition.}
    Its two components can be correlated. A coefficient of
    \(0.5\) therefore does not establish that each component
    explains half the target variance. Final standardization
    also means that the synthetic response has no literal
    millivolt scale.
\end{notes}

\subsection{Preparing the Task for Transformer Training}
\label{sec:synthetic-task-handoff}

\begin{notes}
    \item \textbf{Optional composition descriptors.}
    An optional Magpie-style representation appends ten
    deterministic composition descriptors to the raw table
    before column standardization, giving \(31\) transformer
    inputs. The response construction still uses the original
    \(21\) physical columns. Because standardization is
    column-wise, those first \(21\) standardized inputs remain
    identical to the SCM inputs.
    This option is disabled in the reported search and
    Trial~56 model.

    \item \textbf{Check that the task is usable.}
    The generator checks the generated inputs and targets.
    Unusable draws are rejected and regenerated.
    Missing calibrated-rule scores are a configuration error
    in this mode; the generator requires them to construct
    the informed response.

    \item \textbf{Use the synthetic context/query split.}
    The generated standardized inputs and responses form \(D\).
    Its context and query portions are \(D_{\mathrm{ctx}}\)
    and \(D_{\mathrm{qry}}\), as defined in
    Section~\ref{sec:pfn-pretraining}.
    The transformer receives the context features and labels,
    plus the query features. Query labels are used to calculate
    the training loss.

    \item \textbf{What changes through repeated generation.}
    New tasks can use different compositions, environments,
    MLP/tree mechanisms, analytical rules, and noise.
    Transformer training updates the shared parameters
    \(\phi\) across those tasks. This connects the generator
    described here to the PFN training procedure in
    Section~\ref{sec:pfn-pretraining}.
\end{notes}

\section{Calibrating and Validating the Analytical EPIT Rules}
\label{sec:rule-calibration}

\subsection{Purpose, Inputs, and Eligible Rows}
\label{sec:calibration-scope}

\begin{notes}
    \item \textbf{Where this step belongs.}
    Calibration takes place after the real-data splits are
    frozen and before transformer pretraining begins.
    It prepares the analytical-rule inputs required by the
    generator in Section~\ref{sec:informed-prior}.

    \item \textbf{Two outputs are needed.}
    Each rule needs a validation score to determine how often
    it will be sampled, and a fitted coefficient vector to
    determine how it combines its terms.
    These outputs come from different fits, as described below.

    \item \textbf{Use the same eligible rows for all seven candidates.}
    The candidate formulas in
    Section~\ref{sec:analytical-rule-families} use the
    \(452\) Fe-alloy and Ni--Cr--Mo rows within the development
    partition. The other development rows are outside this
    calibration scope. The saved composition groups and fold
    assignments remain unchanged.

    \item \textbf{Final-test labels are excluded.}
    The calibration script checks the saved source and split
    hashes, then masks final-test EPIT values before fitting
    or scoring any rule. Its calculations use development
    labels only. The broader feature-profile provenance is
    documented separately in
    Section~\ref{sec:feature-profile-provenance}.

    \item \textbf{Candidate-set history.}
    Earlier Al-specific and high-entropy-alloy-specific
    proposals were not retained because their eligible subsets
    were too small or their direct results were unreliable.
    Two historical PREN rules remain as comparison references;
    the synthetic sampler uses the seven candidates defined
    in Section~\ref{sec:analytical-rule-families}.
\end{notes}

\subsection{Preparing One Development Fold}
\label{sec:calibration-fold-inputs}

\begin{notes}
    \item \textbf{Separate context from validation.}
    For fold \(f\), keep its eligible rows as the validation
    queries \(D_{\mathrm{qry},f}\). The eligible rows in the
    other four folds form the labeled context
    \(D_{\mathrm{ctx},f}\).
    Here, the presentation's context/query notation describes
    a real-data rule evaluation.

    \item \textbf{Row counts after the alloy restriction.}
    Folds \(1\)--\(5\) contain
    \((95,96,92,95,74)\) eligible validation rows.
    Their corresponding context sizes are
    \((357,356,360,357,378)\).
    These counts are smaller than the complete-fold counts
    in Section~\ref{sec:dataset-partitions} because calibration
    uses only the two eligible alloy families.

    \item \textbf{Fit preprocessing on the context.}
    Estimate missing-value imputation means for the numeric
    columns required by the rule. Then fit the means and
    scales used by its intermediate signals and final terms.
    All of these estimates use the eligible context rows.
    Apply the saved values unchanged to the validation rows.

    \item \textbf{Construct the term matrices.}
    \(B_{k,\mathrm{ctx},f}\) contains the signed, standardized
    terms of rule \(k\) for the context rows.
    \(B_{k,\mathrm{qry},f}\) contains those terms for the
    validation rows, using the same fitted preprocessing.
    Columns follow the term order listed in
    Section~\ref{sec:analytical-rule-families}.

    \item \textbf{Use target ranks for coefficient fitting.}
    Let \(\mathcal{Z}_{\mathrm{ctx},f}\) denote standardization
    with the relevant mean and scale fitted on the context.
    The fitting target is
    \[
        r_{y,f}
        =
        \mathcal{Z}_{\mathrm{ctx},f}
        \!\left(\operatorname{rank}(y_{\mathrm{ctx},f})\right).
    \]
    Tied EPIT values receive average ranks.
    The objective therefore emphasizes the ordering of EPIT
    values rather than their absolute differences in millivolts.
\end{notes}

\subsection{Preparing the Method-Aware Correction}
\label{sec:method-calibration}

\begin{notes}
    \item \textbf{Group real methods into four broad families.}
    Direct calibration uses potentiodynamic, potentiostatic,
    scratch, and other. These families are used to estimate
    the analytical correction; they do not replace the model's
    original test-method input.
    Scan rate is excluded because it is outside the fixed
    \(21\)-feature schema.

    \item \textbf{Start from an anchored synergy score.}
    Use the context-fitted synergy terms with the predefined
    anchor \(a_{\mathrm{syn}}\), specified in
    Section~\ref{sec:calibration-anchors}:
    \[
        v_f=B_{\mathrm{syn},\mathrm{ctx},f}a_{\mathrm{syn}},
        \qquad
        e_f=r_{y,f}-\mathcal{Z}_{\mathrm{ctx},f}(v_f).
    \]
    The residual \(e_f\) measures what remains unexplained by
    that anchored score. The EPIT target is rank-transformed;
    the anchored prediction is standardized without a rank
    transformation.

    \item \textbf{Estimate and shrink each method-family offset.}
    For family \(j\), let \(n_j\) be its context-row count
    and \(\overline e_j\) its mean residual. Set
    \[
        \widetilde o_j
        =
        \frac{n_j}{n_j+20}\,\overline e_j.
    \]
    Small families receive stronger shrinkage toward zero,
    reducing the influence of offsets estimated from few rows.

    \item \textbf{Center the offsets across context rows.}
    With \(n_f\) context rows in total,
    \[
        \widehat o_j
        =
        \widetilde o_j
        -
        \sum_{\ell}\frac{n_\ell}{n_f}\widetilde o_\ell.
    \]
    Each row receives its family's centered offset. These
    row-wise values are standardized using the context and
    become the fifth column of the method-aware term matrix.

    \item \textbf{Reuse the correction for validation.}
    Apply the fitted family offsets and method-term scaling
    to validation rows. An unseen family receives an offset
    of zero before that saved scaling is applied.
    Validation EPIT values do not enter any of these steps.

    \item \textbf{Then fit the five outer coefficients.}
    Once the correction column is ready, fit the method-aware
    coefficient vector using the same constrained objective as
    the other rules. The base score used to estimate offsets
    uses the predefined synergy anchor, so it does not require
    an earlier fitted synergy coefficient vector.

    \item \textbf{What transfers to synthetic generation.}
    The fitted weight of the method term transfers with the
    other coefficients. The real-data family offsets do not:
    synthetic tasks draw fresh offsets over the \(52\) encoded
    categories, using the construction in
    Section~\ref{sec:analytical-environment-terms}.
\end{notes}

\subsection{Fitting the Relative Coefficients}
\label{sec:coefficient-objective}

\begin{notes}
    \item \textbf{Calibration objective.}
    For candidate \(k\) and fold \(f\), solve
    \begin{equation}
    \begin{aligned}
        \widehat c_{k,f}
        =\underset{c}{\arg\min}\ \bigg\{&
            \frac{1}{n_f}
            \left\|B_{k,\mathrm{ctx},f}c-r_{y,f}\right\|_2^2\\
            &+\gamma\left\|c-a_k\right\|_2^2
        \bigg\},\\
        &\sum_j c_j=1,\qquad 0\leq c_j\leq u_{k,j}.
    \end{aligned}
    \label{eq:rule-calibration}
    \end{equation}
    Here, \(n_f\) is the number of eligible context rows,
    \(a_k\) is a predefined coefficient anchor, and
    \(u_{k,j}\) is the upper bound for coefficient \(j\).

    \item \textbf{Fit term.}
    The first term measures how closely the weighted rule
    scores match the standardized target ranks.
    This is a least-squares objective on ranks.
    It does not directly maximize Spearman correlation;
    Spearman is used afterward to validate the fitted rule.

    \item \textbf{Regularization term.}
    The second term penalizes movement away from the anchor.
    This encourages stable coefficients when several rule
    terms can explain similar patterns.
    Every reported fold fit and final-development fit uses
    \(\gamma=0.1\).

    \item \textbf{Coefficient constraints.}
    Nonnegative coefficients preserve the signs already
    assigned to the rule terms. Requiring their sum to be one
    makes them relative weights. Individual upper bounds
    restrict selected contributions, such as the acidic-pH
    and method terms.

    \item \textbf{What is estimated.}
    The fitted values are dimensionless outer weights.
    The internal constants of the analytical formulas,
    including the chromium thresholds, remain fixed.
    The transformer parameters \(\phi\) are not optimized
    during this calibration stage.

    \item \textbf{Numerical solution.}
    The implementation uses constrained optimization with
    SLSQP, initialized at the anchor. A failed optimization
    stops calibration instead of supplying an invalid
    coefficient vector to later training.
\end{notes}

\subsection{Predefined Anchors and Coefficient Bounds}
\label{sec:calibration-anchors}

\begin{notes}
    \item \textbf{Reading the vectors.}
    Each entry follows the corresponding rule-term order in
    Section~\ref{sec:analytical-rule-families}.
    The anchors below are fixed reference weights used during
    fitting. The fitted weights are reported later in this section.

    \item \textbf{Linear PREN.}
    \[
        a_{\mathrm{PREN}}
        =\frac{(0.575,0.50,0.775)}{1.85},
        \qquad
        u_{\mathrm{PREN}}=(1,1,1).
    \]

    \item \textbf{Synergy and threshold-and-saturation rules.}
    Both use
    \[
        a_{\mathrm{syn}}
        =\frac{(0.575,0.50,0.775,0.20)}{2.05},
        \qquad
        u=(1,1,1,1).
    \]

    \item \textbf{Improved-environment and Fe/Ni-threshold rules.}
    Both use
    \[
        a=(0.42,0.24,0.14,0.14,0.06),
        \qquad
        u=(1,0.65,0.45,0.45,0.15).
    \]

    \item \textbf{Coupled-breakdown rule.}
    \[
        a=(0.55,0.40,0.05),
        \qquad
        u=(1,0.80,0.15).
    \]

    \item \textbf{Method-aware rule.}
    \[
        a=(0.90a_{\mathrm{syn}},0.10),
        \qquad
        u=(1,1,1,1,0.25).
    \]
    Thus, the method term is capped at \(0.25\).
    The acidic-pH term is capped at \(0.15\) in the three
    candidate rules that contain it. Every anchor sums to one.
\end{notes}


\subsection{Validating Each Rule on Unused Fold Rows}
\label{sec:direct-rule-validation}

\begin{notes}
    \item \textbf{Apply the fitted rule to the validation queries.}
    Once preprocessing and coefficients are fixed for fold \(f\),
    calculate
    \[
        \widehat v_{k,f}
        =B_{k,\mathrm{qry},f}\widehat c_{k,f}.
    \]
    These are analytical rule scores. The validation EPIT
    values are used to assess them, not to fit that fold's rule.

    \item \textbf{Measure ordering with Spearman correlation.}
    The score for candidate \(k\) on fold \(f\) is
    \[
        r_{k,f}
        =
        \operatorname{Spearman}
        \!\left(\widehat v_{k,f},y_{\mathrm{qry},f}\right).
    \]
    Repeat the preparation, fitting, and validation for all
    five folds. Each eligible row receives a prediction from
    a rule fitted without that row's validation fold.

    \item \textbf{Average the five fold scores.}
    The saved rule-quality score is
    \begin{equation}
        s_k=\frac{1}{5}\sum_{f=1}^{5}r_{k,f}.
        \label{eq:rule-quality}
    \end{equation}
    Each fold has equal weight, despite their different
    eligible-row counts. This is the mean of five correlations,
    rather than one correlation calculated after pooling
    every validation prediction.

    \item \textbf{Convert quality scores into sampling probabilities.}
    Normalize the seven positive \(s_k\) values using the
    formula in Section~\ref{sec:analytical-target-selection}.
    This gives the probabilities \(p_k\) used by the synthetic
    generator. All seven rules remain available, with better
    direct scores giving a larger sampling probability.
\end{notes}

\subsection{Refitting Once for Synthetic Generation}
\label{sec:final-rule-refit}

\begin{notes}
    \item \textbf{Fit again on all eligible development rows.}
    After fold validation, fit each candidate once more on
    all \(452\) eligible development rows. Refit its
    imputation, scaling, and any method offsets, then solve
    the same coefficient objective with the same anchor,
    bounds, and \(\gamma=0.1\).

    \item \textbf{Save the resulting coefficient vector.}
    This fit supplies the \(c_k\) used during synthetic
    generation. It is a new fit on the complete eligible
    development subset, not an average of the five
    fold-specific coefficient vectors.

    \item \textbf{Keep the validation score from the fold procedure.}
    The final refit's score on its own fitting rows is only
    an in-sample diagnostic. The five-fold \(s_k\) remains
    the quantity used to calculate \(p_k\).

    \item \textbf{Transfer the coefficients to newly generated tasks.}
    Synthetic rule terms are standardized within each generated
    task as described in Section~\ref{sec:informed-prior}.
    The saved coefficients weight those new terms.
    The real-data imputation and scaling states are retained
    for direct-rule evaluation and reproducibility; synthetic
    tasks use their own feature-generation and scaling procedure.
\end{notes}

\subsection{Frozen Scores, Probabilities, and Coefficients}
\label{sec:frozen-rule-results}

\begin{notes}
    \item \textbf{How to read the values.}
    In each entry below, \(s_k\) is the mean direct-validation
    Spearman correlation, \(p_k\) is the synthetic sampling
    probability, and \(c_k\) is the final-development
    coefficient vector. Coefficient entries follow the term
    order in Section~\ref{sec:analytical-rule-families}.
    Values are rounded to four decimal places.

    \item \textbf{Linear PREN.}
    \[
        s_k=0.5174,\qquad p_k=0.1342,\qquad
        c_k=(0.5560,0.4013,0.0427).
    \]

    \item \textbf{Cr--Mo/W synergy.}
    \[
    \begin{aligned}
        s_k&=0.5583,\qquad p_k=0.1448,\\
        c_k&=(0.4331,0.0710,0.2072,0.2887).
    \end{aligned}
    \]

    \item \textbf{Threshold and saturation.}
    \[
    \begin{aligned}
        s_k&=0.5181,\qquad p_k=0.1343,\\
        c_k&=(0.4186,0.1139,0.1960,0.2715).
    \end{aligned}
    \]

    \item \textbf{Improved environment.}
    \[
    \begin{aligned}
        s_k&=0.5850,\qquad p_k=0.1517,\\
        c_k&=(0.5394,0.2414,0.0798,0.1393,0.0000).
    \end{aligned}
    \]

    \item \textbf{Coupled breakdown.}
    \[
        s_k=0.5525,\qquad p_k=0.1433,\qquad
        c_k=(0.4265,0.5062,0.0673).
    \]

    \item \textbf{Fe/Ni chromium threshold.}
    \[
    \begin{aligned}
        s_k&=0.5252,\qquad p_k=0.1362,\\
        c_k&=(0.5233,0.2721,0.1058,0.0988,0.0000).
    \end{aligned}
    \]

    \item \textbf{Method-aware.}
    \[
    \begin{aligned}
        s_k&=0.5999,\qquad p_k=0.1556,\\
        c_k&=(0.3051,0.0000,0.3124,0.2099,0.1727).
    \end{aligned}
    \]
    This candidate has the highest mean direct-validation
    score among the seven and therefore the largest
    sampling probability.

    \item \textbf{Terms assigned zero weight.}
    The acidic-pH term is inactive in the final
    improved-environment and Fe/Ni-threshold fits.
    The combined environment term is inactive in the
    method-aware fit, although its material--chloride and
    temperature--chloride terms remain active.
    These fitted zeros do not establish that the corresponding
    physical effects are absent in real corrosion.

    \item \textbf{Historical PREN references.}
    The older reference uses the material score
    \[
        R_{\mathrm{historical}}
        =\mathrm{Cr}+0.25\mathrm{Ni}
         +3.3\mathrm{Mo}+1.65\mathrm{W}.
    \]
    Its mean fold Spearman score is \(0.4598\) when calibrated
    and evaluated across all \(608\) development rows, and
    \(0.5123\) on the same \(452\)-row Fe/Ni scope used by the
    seven candidates. The \(0.5123\) value is the comparison
    with a matched row population. Neither historical
    reference enters the seven-rule sampler.

    \item \textbf{What these results measure.}
    These correlations assess the analytical formulas directly.
    They are not transformer prediction results.
    A rule's direct score establishes neither the performance
    of a PFN trained with that rule nor its real-data uncertainty
    calibration.
\end{notes}

\subsection{Saved Outputs and the Next Workflow Step}
\label{sec:calibration-outputs}

\begin{notes}
    \item \textbf{Saved calibration records.}
    The directory
    \path{corrosion_datasets/analysis/epit_pipeline/target_rules_v2/}
    contains one JSON artifact per rule,
    \texttt{calibration\_summary.json},
    \texttt{direct\_evaluation\_predictions.csv}, and the
    direct-evaluation HTML report. The records include fold
    scores, fitted coefficients, preprocessing states, and
    the final-development fits.

    \item \textbf{What is passed to transformer training.}
    Training receives all seven rule scores and all
    \(29\) final-development coefficient values.
    Saved floating-point values are preserved when constructing
    the training arguments; the rounded values above are
    only for reading. Artifact hashes bind these values to
    the source data and frozen split.

    \item \textbf{Freeze the calibrated rule distribution.}
    The same scores and coefficient vectors are reused during
    the Optuna search and final transformer training.
    At this point, all eligible development labels have
    influenced the calibrated prior. The next stage searches
    the generator settings and evaluates the resulting
    transformers while keeping these calibration outputs fixed.

    \item \textbf{Implementation references.}
    \path{scripts/epit_pipeline/target_rules.py} constructs
    the real-data rule terms and preprocessing states.
    \path{scripts/epit_pipeline/calibrate_target_rules.py}
    performs coefficient fitting, fold validation, and
    final-development refitting.
\end{notes}

\section{Selecting the Prior and Training the Final Transformer}
\label{sec:model-selection}

\subsection{The Two Selection Decisions}
\label{sec:selection-sequence}

\begin{notes}
    \item \textbf{First choose the generator settings.}
    Starting from the calibrated rules in
    Section~\ref{sec:rule-calibration}, Optuna compares
    candidate settings for the synthetic-data generator.
    One \emph{trial} means one candidate configuration,
    its short transformer-training run, and its development
    evaluation.

    \item \textbf{Then choose a checkpoint.}
    Train a new transformer from scratch using the selected
    configuration for a longer run. Compare checkpoints
    from that run and keep the one with the best development
    score. Choosing the prior configuration and choosing the
    training step are separate decisions.

    \item \textbf{Use development data for both decisions.}
    All trial and checkpoint comparisons use the five frozen
    development folds. The \(152\) final-test rows are
    excluded from these evaluations.

    \item \textbf{Study identity.}
    The reported search is
    \path{epit_pipeline_optuna_empirical_features_scm_target_v7}.
    It uses the corrected generator described in
    Section~\ref{sec:informed-prior}. Earlier studies with
    different target construction or preprocessing are kept
    separate, so their scores do not enter this search.
\end{notes}

\subsection{The Four Searched Parameters}
\label{sec:prior-search-space}

\begin{notes}
    \item \textbf{Informed-task probability.}
    \[
        \rho\in\{0.25,0.50,0.75,1.00\}.
    \]
    This chooses how frequently training uses an informed
    task. The remaining tasks use the generic branch.

    \item \textbf{MLP probability within the informed branch.}
    \[
        p_{\mathrm{MLP}}\in\{0,0.25,0.50,0.70,0.75,1.00\}.
    \]
    This chooses between MLP and tree SCM target heads for
    informed tasks. The generic branch retains its fixed
    MLP/tree probabilities of \(0.7/0.3\).

    \item \textbf{SCM--EPIT response mixture.}
    \[
        \lambda\in[0,1]\qquad\text{(continuous)}.
    \]
    This controls the analytical EPIT contribution in the
    target mixture from Equation~\eqref{eq:scm-epit-target}.

    \item \textbf{Composition perturbation strength.}
    \[
        \tau\in[0,0.15]\qquad\text{(continuous)}.
    \]
    This controls the multiplicative changes around sampled
    alloy templates, as in
    Equation~\eqref{eq:composition-perturbation}.

    \item \textbf{What stays fixed across trials.}
    The empirical profiles, Fe/Ni template-frequency mixture,
    \(21\)-column physical schema, preprocessing policy,
    calibrated rule scores, and calibrated coefficients are
    fixed. Informed tasks always use the physical profile;
    feature permutation, extra block coupling, Dirichlet
    redraw, and Magpie augmentation are disabled.
    The model architecture and optimization settings also
    remain fixed.
\end{notes}

\subsection{How Optuna Proposes and Records Trials}
\label{sec:optuna-trial-proposals}

\begin{notes}
    \item \textbf{Sampler.}
    The search uses Optuna's tree-structured Parzen estimator
    (TPE). It uses earlier trial outcomes to propose settings
    for later trials.

    \item \textbf{Initial random sampling.}
    The setting is \texttt{n\_startup\_trials=10}.
    Sampling remains random while fewer than ten completed
    or pruned trials are available. With concurrent workers,
    more than ten trials can start randomly before enough
    results have finished.

    \item \textbf{Shared search.}
    Workers contribute to the same Optuna journal.
    Each worker uses one GPU and one training process.
    The records include trial parameters, commands, scores,
    seeds, worker settings, and hashes of the source,
    split, and calibration artifacts.

    \item \textbf{Separate sampler and training seeds.}
    The NumPy and Torch training seeds are both \(42\).
    Optuna sampler seeds are recorded per worker and trial;
    the selected Trial~56 records sampler seed \(66373\).
    The sampler seed controls proposed configurations, while
    the training seeds control randomness within the
    generated tasks and model-training run.
\end{notes}

\subsection{Short Transformer Training for Each Trial}
\label{sec:proxy-training}

\begin{notes}
    \item \textbf{Train one model from scratch.}
    Each trial trains a transformer for \(1{,}000\) steps
    using tasks from its candidate generator configuration.
    A step is a parameter update from a batch of synthetic
    tasks. The resulting checkpoint supplies all five
    development-fold evaluations for that trial.

    \item \textbf{Fixed Stage~1 setup.}
    Training uses synthetic tables of \(1{,}024\) rows.
    The configured synthetic context fraction ranges from
    \(0.1\) to \(0.9\).
    Each update uses a batch of \(512\) tasks, processed
    in microbatches of four.
    The regression loss is CRPS, as described in
    Section~\ref{sec:tabicl-regression}.

    \item \textbf{Fixed optimizer and schedule.}
    The local trainer uses AdamW with learning rate
    \(10^{-4}\), a cosine schedule with warmup, and gradient
    clipping at \(1.0\).
    The warmup proportion is \(0.02\) of the
    \(10{,}000\)-step scheduler horizon, giving \(200\)
    warmup steps.

    \item \textbf{The short run uses the longer schedule.}
    The proxy stops at step \(1{,}000\), but the scheduler
    still has a \(10{,}000\)-step horizon.
    It therefore evaluates an early part of the intended
    training schedule. The schedule is not compressed to
    finish its decay within the proxy run.

    \item \textbf{Data generation and loading.}
    Synthetic tasks are generated during training on CPU.
    The recorded configuration uses one parallel
    prior-generation job, four data-loader workers, and
    a prefetch factor of four. Model optimization runs on
    the GPU.

    \item \textbf{Purpose of the proxy.}
    Short runs make it possible to compare more generator
    configurations before committing to longer training.
    Their ranking is an estimate of which configurations are
    useful; it does not guarantee the same ranking after
    every possible training duration.
\end{notes}

\subsection{Evaluating a Trial on the Development Folds}
\label{sec:transformer-development-evaluation}

\begin{notes}
    \item \textbf{Use the complete development population.}
    Transformer evaluation covers all \(608\) development
    rows, across all alloy classes.
    Each fold supplies \(486\)--\(487\) labeled context rows
    and \(121\)--\(122\) validation query rows.
    The rule-calibration restriction to \(452\) Fe/Ni rows
    does not apply here.

    \item \textbf{Reuse the trained checkpoint.}
    For each fold, supply its context features and EPIT
    labels together with its query features.
    Apply the context-fitted real-data preprocessing from
    Section~\ref{sec:real-data-representation}.
    The context changes between folds while the learned
    transformer parameters \(\phi\) remain fixed.

    \item \textbf{Calculate the trial objective.}
    Let \(t\) identify a trial and
    \(\widehat y_{t,f}\) its EPIT predictions on fold \(f\).
    Optuna maximizes
    \begin{equation}
        S_t
        =
        \frac{1}{5}\sum_{f=1}^{5}
        \operatorname{Spearman}
        \!\left(\widehat y_{t,f},y_{\mathrm{qry},f}\right).
        \label{eq:optuna-objective}
    \end{equation}
    This is an equal-weight mean of the five fold correlations.

    \item \textbf{Why ranking is the primary criterion.}
    EPIT spans a wide numerical range and includes extreme
    observations. Spearman measures agreement in the ordering
    of predictions and observations, while magnitude-based
    errors can be strongly affected by a few large errors.

    \item \textbf{Keep magnitude-based diagnostics.}
    MAE, RMSE, \(R^2\), and variation across folds are also
    recorded. They help show whether good ordering is
    accompanied by accurate millivolt predictions.
    They do not determine the winning trial.

    \item \textbf{Read saved metric names in context.}
    Fields such as \texttt{mean\_test\_spearman} in these
    evaluator outputs refer to the held-out development
    folds. They are development scores, despite the word
    \texttt{test} in the field name.
\end{notes}


\subsection{One Fixed Inference Configuration}
\label{sec:selection-inference-settings}

\begin{notes}
    \item \textbf{Use the same prediction settings throughout.}
    Proxy-fold evaluation, final-run checkpoint selection,
    and the locked final inference configuration use the
    same feature representation and prediction settings.
    This keeps changes in evaluation settings from affecting
    which trial or checkpoint wins.

    \item \textbf{Requested and effective ensemble sizes.}
    The estimator is configured with
    \texttt{n\_estimators=8},
    \texttt{feat\_shuffle\_method="none"}, and
    \texttt{norm\_methods=["none"]}.
    These options produce one feature order and one numerical
    view, giving one actual inference configuration.
    Eight is the requested upper limit; the effective
    ensemble size here is one.

    \item \textbf{Meaning of the normalization setting.}
    The \texttt{none} option disables additional
    power-transformed inference views. The estimator's
    ordinary numeric preprocessing still applies.

    \item \textbf{Outputs and row use.}
    Predictions use the median regression output, transformed
    back to millivolts, with regression-uncertainty output
    disabled. Inference uses random state \(42\).
    The locked \texttt{max\_samples\_per\_task=0} setting
    means all available rows are used.
    The reported representation has \(21\) inputs and no
    Magpie descriptors.
\end{notes}

\subsection{How to Interpret the Development Scores}
\label{sec:development-score-interpretation}

\begin{notes}
    \item \textbf{The analytical prior stays fixed within the search.}
    Before Optuna begins, rule probabilities use results from
    all five calibration folds, and rule coefficients are
    fitted on all \(452\) eligible development rows.
    These values are reused for every transformer fold;
    calibration is not repeated inside each fold evaluation.

    \item \textbf{Validation labels have already influenced the prior.}
    The eligible validation rows are absent from the labeled
    context for their transformer fold, but their EPIT values
    have contributed to the frozen calibration outputs.
    This is why the transformer evaluation is not a complete
    cross-validation of the whole calibration-and-training
    procedure.

    \item \textbf{Use these scores for selection.}
    The development scores guide both configuration selection
    and checkpoint selection. They are not unbiased estimates
    of performance on new data.
    Final-test labels are excluded from these decisions;
    the feature-profile boundary remains as described in
    Section~\ref{sec:feature-profile-provenance}.

    \item \textbf{Direct-rule and transformer scores have different scopes.}
    The direct scores in
    Section~\ref{sec:frozen-rule-results} cover the
    \(452\)-row Fe/Ni subset. The transformer scores cover
    all \(608\) development rows.
    Comparing their numerical values alone does not measure
    the improvement from replacing a formula with a transformer.
\end{notes}

\subsection{The Selected Optuna Configuration}
\label{sec:selected-prior-configuration}

\begin{notes}
    \item \textbf{Selection uses the available completed trials.}
    A candidate must be completed and have a finite
    development objective. The final run takes the completed
    candidate with the highest \(S_t\) available at selection
    time. The recorded decision remains fixed even if the
    study later receives more results.

    \item \textbf{Recorded study state.}
    At selection, the study contained \(78\) trials, with
    six still marked as running.
    Trial~56 was explicitly selected as the best completed
    trial then available. The unfinished trial identifiers
    and states are retained in the manifest.

    \item \textbf{Selected parameter values.}
    Rounded to four decimal places where needed,
    \[
        \rho=0.75,\qquad
        p_{\mathrm{MLP}}=0.50,\qquad
        \lambda=0.2048,\qquad
        \tau=0.1028.
    \]
    Thus, \(75\%\) of tasks use the informed branch,
    and its MLP and tree heads are equally likely.
    The informed response mixture has coefficients of
    approximately \(0.7952\) for the standardized SCM
    response and \(0.2048\) for the standardized analytical
    response.

    \item \textbf{Proxy development results.}
    At the proxy checkpoint, step \(1{,}000\), Trial~56
    has the following means across the five development folds:
    \[
    \begin{aligned}
        \text{Spearman}&=0.6903,&
        \text{MAE}&=211.0\ \text{mV},\\
        \text{RMSE}&=298.0\ \text{mV},&
        R^2&=0.4373.
    \end{aligned}
    \]
    The Spearman value determines its selection.
\end{notes}

\subsection{Long Training and Checkpoint Selection}
\label{sec:final-training-and-selection}

\begin{notes}
    \item \textbf{Start a new training run.}
    Train the selected configuration from scratch for
    \(10{,}000\) Stage~1 steps, retaining the fixed
    optimization settings and calibrated rule distribution.
    Transformer weight updates use synthetic tasks.
    Real development labels are used through prior
    calibration and as context and scoring targets during
    development evaluation.

    \item \textbf{Save candidate checkpoints.}
    Permanent checkpoints are saved every \(500\) steps.
    After the training run finishes, evaluate the \(20\)
    saved checkpoints at steps
    \(500,1000,\ldots,10000\) on the five development folds.

    \item \textbf{Select by the same primary metric.}
    Let \(S_{\mathrm{dev}}(h)\) be the mean development-fold
    Spearman correlation for the checkpoint at training step
    \(h\). Select
    \begin{equation}
        h_{\mathrm{selected}}
        =
        \underset{h\in\{500,1000,\ldots,10000\}}{\arg\max}
        \ S_{\mathrm{dev}}(h).
        \label{eq:checkpoint-selection}
    \end{equation}
    All candidates use the inference settings from
    Section~\ref{sec:selection-inference-settings}.
    MAE, RMSE, and \(R^2\) remain diagnostics.

    \item \textbf{Selected checkpoint and development results.}
    Step \(2{,}500\) has the highest mean development
    Spearman correlation:
    \[
    \begin{aligned}
        \text{Spearman}&=0.6966,&
        \text{MAE}&=191.3\ \text{mV},\\
        \text{RMSE}&=280.9\ \text{mV},&
        R^2&=0.4717.
    \end{aligned}
    \]
    The \(10{,}000\)-step run was completed, but development
    evaluation selected these earlier saved weights.
    Final-test scores do not enter this choice.

    \item \textbf{Training-repeat scope.}
    The reported result uses one final training run with
    the recorded seeds. Variation across development folds
    describes different context/query partitions; it does
    not measure variation across independently trained
    final models.
\end{notes}

\subsection{Freezing the Model for Final Evaluation}
\label{sec:final-model-lock}

\begin{notes}
    \item \textbf{Save the complete selection record.}
    The final-model manifest records Trial~56, the effective
    training configuration, rule artifacts, split hashes,
    all checkpoint evaluations, the selected checkpoint
    hash, and the final inference settings.
    A separate lock binds the manifest to the selected
    checkpoint file.

    \item \textbf{Recorded run location.}
    The manifest and lock are stored under
    \path{corrosion_datasets/analysis/epit_pipeline/final_v1/epit_pipeline_optuna_empirical_features_scm_target_v7/}
    as \texttt{final\_model\_manifest.json} and
    \texttt{final\_model\_lock.json}.
    They identify \texttt{step-2500.ckpt} as the model to
    carry into final evaluation.

    \item \textbf{Keep changes identifiable.}
    Existing final manifests and locks are not overwritten.
    A run with parameter overrides receives a distinct
    identity and records those overrides.
    The reported Trial~56 final run has no such overrides.

    \item \textbf{Workflow entry points.}
    \path{scripts/epit_pipeline/run_optuna.py} manages the
    search, and
    \path{scripts/epit_pipeline/evaluate_optuna_folds.py}
    evaluates development folds.
    \path{scripts/epit_pipeline/train_final.py} performs
    the long training run, checkpoint comparison, and
    final-model freezing.
\end{notes}

\section{Final Evaluation and Results}
\label{sec:final-evaluation}

\subsection{From the Frozen Model to Final-Test Predictions}
\label{sec:final-prediction-task}

\begin{notes}
    \item \textbf{Start with the selected model.}
    CorrPFN is Trial~56 at step \(2{,}500\), selected using
    development scores and locked as described in
    Section~\ref{sec:final-model-lock}.
    The final-test scores below do not determine this choice.

    \item \textbf{Supply the final context and queries.}
    All \(608\) development rows, including their EPIT
    labels, form \(D_{\mathrm{ctx}}\).
    The \(152\) final-test rows supply query features.
    These queries include all alloy classes assigned to the
    final-test partition.

    \item \textbf{Predict with fixed transformer weights.}
    The trained model uses the labeled context to predict
    query EPIT values. The analytical corrosion rules have
    already shaped its synthetic pretraining; they are not
    evaluated as an extra correction during final inference.

    \item \textbf{Keep the input and output definitions fixed.}
    The comparison uses the \(21\) raw input columns without
    Magpie augmentation. TabICL preprocessing is fitted on
    the development context, and its point predictions are
    medians returned on the millivolt scale.
    The locked inference configuration is documented in
    Section~\ref{sec:selection-inference-settings};
    the retained execution records are discussed in
    Section~\ref{sec:final-result-provenance}.
\end{notes}

\subsection{Comparison Models and Their Selection}
\label{sec:comparison-models}

\begin{notes}
    \item \textbf{Local generic baseline.}
    This model was trained on the generic
    \texttt{mix\_scm} prior, with MLP/tree probabilities of
    \(0.7/0.3\). It provides a local reference with the same
    transformer architecture as CorrPFN.

    \item \textbf{What is matched in the saved baseline records.}
    The model configurations and parameter-tensor shapes
    match. At step \(2{,}500\), the recorded optimizer
    hyperparameters and learning-rate settings also match.
    The available checkpoint metadata does not establish
    equality of every training setting.

    \item \textbf{Generic baseline at the same training step.}
    Step \(2{,}500\) matches the step of the selected CorrPFN
    checkpoint. This comparison holds the recorded update
    count fixed.

    \item \textbf{Generic baseline selected by development scores.}
    A second comparison applies the same selection criterion
    used for CorrPFN: highest mean development-fold Spearman.
    Among \(20\) saved baseline checkpoints, this selects
    step \(7{,}000\), with development Spearman \(0.6723\).
    This comparison was added retrospectively using the
    development evaluations recorded on 2 September 2026;
    it was not frozen in the original CorrPFN selection lock.

    \item \textbf{Pretrained TabICL V2.}
    This reference uses the authors' released
    \texttt{tabicl-regressor-v2-20260212.ckpt} weights.
    Its synthetic generator and pretraining procedure differ
    from the local runs, even though the local architecture
    also contains V2 components.
    The comparison therefore changes more than the
    corrosion-informed part of the prior.

    \item \textbf{CatBoost.}
    This is the conventional supervised regression reference.
    It is fitted on all \(608\) development rows, then
    predicts the same \(152\) final-test rows.
    The recorded CatBoost version is \(1.2.10\).

    \item \textbf{Fixed CatBoost settings.}
    The run uses \(1{,}000\) boosting iterations, depth \(6\),
    learning rate \(0.03\), L2 regularization \(3\), and
    random seed \(42\). Its training objective is RMSE.
    It runs on CPU with eight threads.
    There is no CatBoost hyperparameter search, validation
    set during fitting, or early stopping in this comparison.

    \item \textbf{CatBoost input representation.}
    CatBoost receives the same raw feature columns, handles
    numeric missing values directly, and treats test method
    as categorical. Missing categories are replaced by a
    dedicated string label before fitting.

    \item \textbf{References independent of local checkpoint step.}
    Pretrained V2 and CatBoost do not change when a different
    local checkpoint is inspected. The two reported generic
    baseline rows represent different saved states of the
    same local training run.
\end{notes}

\subsection{Calculating and Reading the Final Metrics}
\label{sec:final-metrics}

\begin{notes}
    \item \textbf{Evaluation quantities.}
    Let \(m=152\) be the number of final-test queries.
    For query \(i\), \(y_i^\ast\) is the observed EPIT value
    and \(\widehat y_i^\ast\) is its prediction.
    All four metrics below are calculated on this same
    final-test partition. They are not the five-fold
    development means from Section~\ref{sec:model-selection}.

    \item \textbf{Spearman correlation: higher is better.}
    Spearman measures agreement between the ranks of
    \(\widehat y^\ast\) and \(y^\ast\).
    It measures ordering without requiring accurate
    millivolt differences between predictions.

    \item \textbf{Mean absolute error: lower is better.}
    \begin{equation}
        \operatorname{MAE}
        =
        \frac{1}{m}\sum_{i=1}^{m}
        \left|\widehat y_i^\ast-y_i^\ast\right|.
        \label{eq:final-mae}
    \end{equation}
    MAE reports the average absolute prediction error in
    millivolts.

    \item \textbf{Root mean squared error: lower is better.}
    \begin{equation}
        \operatorname{RMSE}
        =
        \sqrt{\frac{1}{m}\sum_{i=1}^{m}
        \left(\widehat y_i^\ast-y_i^\ast\right)^2}.
        \label{eq:final-rmse}
    \end{equation}
    RMSE is also in millivolts and gives larger errors more
    influence than MAE.

    \item \textbf{Coefficient of determination: higher is better.}
    With \(\overline y^\ast\) denoting the final-test mean,
    \begin{equation}
        R^2
        =
        1-
        \frac{\sum_{i=1}^{m}
            (\widehat y_i^\ast-y_i^\ast)^2}
             {\sum_{i=1}^{m}
            (y_i^\ast-\overline y^\ast)^2}.
        \label{eq:final-r-squared}
    \end{equation}
    A negative value means that the squared prediction
    error exceeds the error of a constant predictor equal
    to the final-test mean. That mean is the mathematical
    reference in this metric; it is not a separately trained
    comparison model.

    \item \textbf{Scope of the reported outputs.}
    These metrics assess point predictions.
    They do not evaluate the calibration or coverage of
    predictive intervals. Regression-uncertainty output is
    disabled in the recorded comparison.
\end{notes}

\subsection{Results on the 152 Final-Test Rows}
\label{sec:held-out-results}

\begin{notes}
    \item \textbf{Reading the results.}
    Correlations and \(R^2\) are rounded to four decimal
    places; errors are rounded to one decimal place.
    Every model below uses the same final-test rows.

    \item \textbf{CorrPFN, selected step \(2{,}500\).}
    \[
    \begin{aligned}
        \text{Spearman}&=0.7681,&
        \text{MAE}&=186.8\ \text{mV},\\
        \text{RMSE}&=268.7\ \text{mV},&
        R^2&=0.5294.
    \end{aligned}
    \]

    \item \textbf{Pretrained TabICL V2.}
    \[
    \begin{aligned}
        \text{Spearman}&=0.6011,&
        \text{MAE}&=206.8\ \text{mV},\\
        \text{RMSE}&=343.8\ \text{mV},&
        R^2&=0.2294.
    \end{aligned}
    \]

    \item \textbf{CatBoost.}
    \[
    \begin{aligned}
        \text{Spearman}&=0.6090,&
        \text{MAE}&=231.7\ \text{mV},\\
        \text{RMSE}&=323.1\ \text{mV},&
        R^2&=0.3194.
    \end{aligned}
    \]

    \item \textbf{Generic baseline, same step \(2{,}500\).}
    \[
    \begin{aligned}
        \text{Spearman}&=0.2829,&
        \text{MAE}&=326.4\ \text{mV},\\
        \text{RMSE}&=548.8\ \text{mV},&
        R^2&=-0.9630.
    \end{aligned}
    \]

    \item \textbf{Generic baseline, development-selected step \(7{,}000\).}
    This is the additional retrospective comparison described
    in Section~\ref{sec:comparison-models}.
    \[
    \begin{aligned}
        \text{Spearman}&=0.5474,&
        \text{MAE}&=226.4\ \text{mV},\\
        \text{RMSE}&=337.8\ \text{mV},&
        R^2&=0.2561.
    \end{aligned}
    \]
\end{notes}

\subsection{What the Recorded Comparison Shows}
\label{sec:final-result-interpretation}

\begin{notes}
    \item \textbf{Best values across the four reported metrics.}
    CorrPFN has the highest Spearman and \(R^2\), and the
    lowest MAE and RMSE, among these recorded comparisons.
    The improvement therefore includes both ordering and
    millivolt-scale prediction accuracy.

    \item \textbf{Difference from pretrained TabICL V2.}
    Using the unrounded metric values, CorrPFN reduces MAE
    by \(9.7\%\), reduces RMSE by \(21.9\%\), and increases
    Spearman by \(0.1670\).
    The error reductions are relative percentages; the
    correlation difference is an absolute change.

    \item \textbf{The stronger generic comparison matters.}
    The development-selected generic checkpoint performs
    substantially better than its step-\(2{,}500\) checkpoint
    on this final-test partition. CorrPFN still has better
    values on all four metrics when compared with that
    step-\(7{,}000\) reference.

    \item \textbf{Development and final-test scores are different evaluations.}
    CorrPFN's final-test Spearman is higher than its
    development-fold mean, while MAE and RMSE are lower.
    Both the query rows and context size change: final
    inference uses \(608\) context rows instead of
    \(486\)--\(487\).
    These numbers alone cannot isolate the effect of more
    context or establish statistical significance.

    \item \textbf{Scope of the finding.}
    The result supports the complete informed configuration
    on this dataset and partition.
    It does not separately establish the benefit of empirical
    features, rule calibration, generic-task mixing, or any
    individual rule. Those contributions require matched
    ablations of the final system.
\end{notes}

\subsection{Later Checkpoint Diagnostics}
\label{sec:post-selection-diagnostics}

\begin{notes}
    \item \textbf{Keep the chronology explicit.}
    CorrPFN was locked on 1 September 2026.
    The comparison run recorded on 10 September evaluated
    \(25\) common checkpoints of the two local models on
    the final-test set. This included additional saved
    checkpoints beyond the \(20\) candidates used for
    development selection.

    \item \textbf{The best-looking final-test checkpoint was different.}
    CorrPFN step \(2{,}000\) produced the strongest recorded
    final-test values across all four reported metrics.
    This was discovered by inspecting final-test results,
    after the development-based selection.

    \item \textbf{Retain the original selected model.}
    Step \(2{,}500\) remains the reported CorrPFN checkpoint.
    Replacing it with step \(2{,}000\) because of the later
    final-test scores would use the evaluation labels to
    choose the model.

    \item \textbf{Role of the additional results.}
    The other checkpoints provide diagnostics of how
    performance changes through training.
    A new model or configuration chosen using these
    final-test diagnostics would need a new independent
    evaluation to support a fresh performance claim.
\end{notes}

\subsection{Where the Final Results Come From}
\label{sec:final-result-provenance}

\begin{notes}
    \item \textbf{Dedicated final-evaluation entry point.}
    \path{scripts/epit_pipeline/evaluate_final.py}
    checks the frozen model and split, verifies the expected
    feature schema, and predicts with the selected checkpoint.
    A successful run writes
    \texttt{final\_evaluation\_manifest.json} and prevents
    a completed evaluation in that output directory from
    being silently replaced.

    \item \textbf{Retained source of the reported metrics.}
    The scores in Section~\ref{sec:held-out-results} come
    from the separate comparison run:
    \path{corrosion_datasets/analysis/epit_pipeline/trial56_final_test_all_checkpoints_11172/summary.csv}.
    Use its step-\(2{,}500\) entries for CorrPFN, pretrained
    V2, CatBoost, and the equal-step generic baseline, plus
    the generic step-\(7{,}000\) entry.
    The accompanying \texttt{results.json} records the
    comparison date, checkpoint paths, and available settings.

    \item \textbf{Selection record and execution record.}
    The comparison identifies the same CorrPFN checkpoint
    path and step as the final-model manifest.
    That manifest and its lock document the earlier
    development-based selection.
    A completion manifest from the dedicated final-evaluation
    entry point is absent from the retained artifacts;
    the reported metrics are therefore attributed to the
    comparison output above.

    \item \textbf{One missing setting in the comparison metadata.}
    The retained comparison JSON records the estimator limit,
    fixed feature order, and median output, but does not
    record the normalization-view setting.
    It therefore does not independently confirm that the
    locked \texttt{norm\_methods=["none"]} setting, and hence
    the single effective inference view described in
    Section~\ref{sec:selection-inference-settings}, was
    used in that comparison run.
    The missing entry is a limit of the execution record,
    not evidence that a different setting was used.

    \item \textbf{Source of the generic checkpoint selection.}
    The saved development results are
    \path{corrosion_datasets/analysis/epit_pipeline/trial56_checkpoint_folds_66378/summary.csv}.
    The generic step-\(7{,}000\) choice follows the largest
    mean development-fold Spearman in that file.
    Its retrospective status is kept explicit alongside
    the final-test comparison.
\end{notes}

\section{Discussion, Limitations, and Next Experiments}
\label{sec:discussion}

\subsection{What Earlier Development Explains}
\label{sec:development-lessons}

\begin{notes}
    \item \textbf{A good synthetic-data diagnostic need not select a good transformer.}
    Earlier experiments trained simpler models, including
    Ridge regression, on synthetic data and tested their
    predictions on real EPIT data. Their scores did not
    reliably rank prior configurations in the same order as
    downstream transformer performance.
    This explains the final division of work: direct rule
    calibration shapes the analytical contribution, while
    actual transformer evaluations select the searched
    prior settings.

    \item \textbf{The target must remain connected to the observed inputs.}
    Earlier generators could transform a composition after
    generating a target from latent variables. That
    transformation could discard information needed to
    predict the target.
    The final construction generates the physical features
    first, then calculates the response from them, as
    described in Section~\ref{sec:generation-order}.

    \item \textbf{Both the target formula and its inputs matter.}
    The v7 correction restored the calibrated EPIT contribution
    that the preceding SCM-only implementation had bypassed.
    It also preserved raw inputs for the corrosion formulas
    and aligned the standardized inputs supplied to the SCM
    and transformer. These changes make the intended
    construction explicit; they do not, by themselves,
    quantify its predictive benefit.

    \item \textbf{Historical experiments explain decisions.}
    The earlier runs in \texttt{main2.tex} used different
    generators, splits, and training settings.
    They provide development evidence, but they cannot
    substitute for comparisons that change one part of the
    final system while keeping the rest fixed.
\end{notes}

\subsection{How Far the Final-Test Claim Extends}
\label{sec:generalization-boundary}

\begin{notes}
    \item \textbf{The protected boundary concerns the revised workflow.}
    After the split was fixed, final-test EPIT labels were
    excluded from analytical-rule calibration, Optuna
    selection, and checkpoint selection.
    However, the complete \(760\)-row corpus had influenced
    historical prior development before that split existed.
    The final-test set was therefore not a newly obtained
    dataset that had never informed the project.

    \item \textbf{Unseen labeled groups do not mean unseen feature values.}
    As documented in
    Section~\ref{sec:feature-profile-provenance}, the saved
    synthetic profiles include feature information from
    the full corpus. A final-test composition can therefore
    influence pretraining as an unlabeled template, even
    though its composition group is absent from the labeled
    prediction context.
    The result measures generalization under this
    dataset-adapted prior.

    \item \textbf{The split answers a composition-group question.}
    Source publications may occur on both sides of the
    split, and all observations come from the same database.
    The result does not establish performance in a new
    laboratory, publication, alloy database, or measurement
    campaign. Those settings can differ in ways that the
    current split does not separate.

    \item \textbf{The existing final-test set is now part of the experiment history.}
    The later checkpoint diagnostics are described in
    Section~\ref{sec:post-selection-diagnostics}.
    Further changes motivated by those results need fresh
    independent evaluation; retaining the original checkpoint
    does not make the inspected test set unseen again.
\end{notes}

\subsection{Uncertainty in the Model Comparison}
\label{sec:statistical-limitations}

\begin{notes}
    \item \textbf{One final training run does not show stability across runs.}
    Different random seeds can change synthetic tasks and
    transformer training. The current results do not measure
    that variation. Variation across development folds
    concerns different context/query partitions and does
    not replace repeated training.

    \item \textbf{The test rows are not necessarily independent observations.}
    Several measurements can belong to one composition group
    or source study. Consequently, \(152\) test rows do not
    automatically provide the same amount of independent
    evidence as \(152\) unrelated experiments.
    The reported metric differences have no accompanying
    assessment of this sampling uncertainty.

    \item \textbf{Development selection has two sources of uncertainty.}
    The prior is selected after \(1{,}000\) proxy-training
    steps, while final training runs for \(10{,}000\) steps.
    A configuration that learns quickly need not perform
    best after longer training.
    Also, the calibration outputs and development folds are
    reused during selection, as explained in
    Section~\ref{sec:development-score-interpretation}.
    The best development score is therefore a selection
    result, not an unbiased estimate for new data.

    \item \textbf{The comparisons apply to the recorded configurations.}
    The generic model's architecture is matched, but the
    retained metadata does not prove that every training
    setting is identical. Pretrained V2 has a different
    pretraining procedure, and CatBoost uses fixed settings
    without its own search.
    The results do not establish that CorrPFN beats every
    possible configuration of these model classes.

    \item \textbf{Prediction uncertainty is a separate question.}
    Uncertainty about a new EPIT value concerns the model's
    predictive distribution. Uncertainty about the reported
    improvement concerns variation across datasets and
    training runs. Good point-prediction scores alone do
    not establish either calibrated predictive intervals
    or a statistically stable advantage.
\end{notes}

\subsection{Limits of the Synthetic Data and Corrosion Rules}
\label{sec:physical-limitations}

\begin{notes}
    \item \textbf{Composition diversity is limited by the template bank.}
    Perturbation changes amounts of elements already present
    in a template; it does not introduce an absent element.
    The Ni--Cr--Mo branch has only \(17\) distinct templates,
    compared with \(298\) Fe-alloy templates.
    Generating many rows from these templates does not
    create equally broad evidence about new alloy systems.

    \item \textbf{Synthetic coverage differs from evaluation coverage.}
    The informed composition generator and rule calibration
    focus on Fe and Ni--Cr--Mo alloys.
    Other alloy classes still occur in the evaluated task.
    Overall test metrics combine these classes and do not
    establish equally good performance for each class.

    \item \textbf{Plausible compositions are not guaranteed physical materials.}
    Composition closure enforces the total element amount,
    but it does not establish phase stability, microstructure,
    heat-treatment response, or passive-film behavior.
    The generator is not a thermodynamic or mechanistic
    corrosion simulator.

    \item \textbf{Recombining alloys and environments is an assumption.}
    Each sampled environmental tuple preserves its observed
    temperature, chloride, pH, and method combination.
    Composition is drawn independently, so the original
    alloy--experiment pairing is lost.
    The resulting combinations may be useful training tasks,
    but their feasibility and experimental relevance have
    not been established for every draw.

    \item \textbf{The equations express selected tendencies.}
    The SCM adds randomly sampled nonlinear relationships,
    while the analytical rules add chosen corrosion-related
    structure. Neither is a complete law of pitting.
    Rule calibration and synthetic target standardization
    also mean that the generating score is not a physical
    EPIT prediction in millivolts.
    The model obtains its real prediction scale from the
    labeled context through the regression procedure in
    Section~\ref{sec:tabicl-regression}.
\end{notes}

\subsection{Limits of Test-Method and Feature Representation}
\label{sec:representation-limitations}

\begin{notes}
    \item \textbf{Arbitrary method numbers can affect a numeric model.}
    As explained in Section~\ref{sec:real-data-representation},
    method codes identify categories without a physical
    order or distance. Nevertheless, a numeric SCM or
    transformer can respond to their numerical spacing.
    For example, it can treat codes \(1\) and \(2\) as closer
    than codes \(1\) and \(40\), although the encoding gives
    no experimental reason for that relationship.

    \item \textbf{Synthetic method effects are not copied protocol corrections.}
    The method-aware rule uses category-specific offsets.
    Its calibrated term weight transfers to synthetic
    generation, but the fitted real method-family offsets
    do not: new offsets are sampled for synthetic tasks.
    This introduces variation between methods without
    reproducing the estimated effect of each real protocol.

    \item \textbf{The input schema omits some relevant experimental detail.}
    A method category does not fully describe a measurement
    procedure. The \(21\) inputs do not explicitly encode
    every detail of surface preparation, microstructure,
    heat treatment, or scan rate.
    The model cannot directly distinguish two rows using
    information that those rows do not contain.

    \item \textbf{The reported model depends on the documented schema.}
    Fixed feature order assigns each column a specific role.
    Changing units, column positions, or the method
    representation requires compatible preprocessing and
    validation. Additional descriptors or a new categorical
    representation must also be used consistently during
    synthetic pretraining and real-data evaluation.
\end{notes}

\subsection{Next Experiments and the Questions They Would Answer}
\label{sec:next-experiments}

\begin{notes}
    \item \textbf{Status of these proposals.}
    The experiments below are follow-up work, not additional
    evidence already included in the reported results.

    \item \textbf{Use a fresh evaluation with development-only profiles.}
    Rebuild composition templates, environmental records,
    imputation statistics, and method mappings using only
    the new development data. Freeze the entire procedure
    before evaluating independent measurements.
    An external dataset or a split that holds out entire
    source studies would test transfer across data sources.

    \item \textbf{Repeat training and quantify uncertainty.}
    Train CorrPFN and its local generic comparator with
    several predefined seeds under matched settings.
    Select checkpoints using development data for every run.
    Compare models on the same queries and assess metric
    differences while accounting for composition groups
    and shared sources.
    This separates variation from training from variation
    in the measured test sample.

    \item \textbf{Change one part of the final generator at a time.}
    An ablation is a controlled comparison that removes or
    changes one component.
    Setting \(\lambda=0\) tests the contribution of the
    analytical response while retaining informed features;
    setting \(\rho=1\) removes generic-task mixing.
    Calibrated versus predefined coefficients, or
    score-weighted versus uniform rule sampling, answer
    different questions and should be separate comparisons.
    Keep the training budget, feature representation,
    development selection, and evaluation protocol matched.

    \item \textbf{Check whether the short proxy ranks longer runs well.}
    Train several promising prior configurations beyond
    \(1{,}000\) steps and compare their development trajectories.
    This tests whether the current search favors fast early
    learning at the expense of later performance.

    \item \textbf{Test a categorical treatment of method.}
    Compare the current integer column with one-hot encoding
    or a representation learned specifically for categories.
    First, a diagnostic can consistently relabel methods in
    both context and queries to measure sensitivity to
    arbitrary code values.
    A replacement encoding requires compatible synthetic
    training and evaluation; changing only the final-test
    input format is not a complete alternative model.

    \item \textbf{Locate the remaining prediction errors.}
    Report errors by alloy family, source study, method,
    and environmental range, together with subgroup sizes.
    If predictive intervals are needed, also test their
    empirical coverage and width on independent queries.
    These checks would show where an overall good score
    conceals weak coverage or unreliable uncertainty.

    \item \textbf{Save a complete record of the next final evaluation.}
    Use the frozen evaluation entry point and retain its
    completion manifest, predictions, and all effective
    inference settings, including normalization views.
    This addresses the execution-record gap identified in
    Section~\ref{sec:final-result-provenance} and makes the
    next comparison easier to reproduce.
\end{notes}

\section{Optional Magpie-Style Composition Descriptors}
\label{sec:magpie-descriptors}

\subsection{Purpose and Status in This Report}
\label{sec:magpie-purpose}

\begin{notes}
    \item \textbf{This is an optional extension.}
    The reported v7 search and selected CorrPFN model use
    \(21\) inputs with Magpie descriptors disabled.
    This section documents the implemented extension
    introduced briefly in
    Section~\ref{sec:synthetic-task-handoff}.

    \item \textbf{What a composition descriptor is.}
    A descriptor is an additional number calculated from
    the element amounts and a table of elemental properties.
    For example, it can summarize how different the
    constituent elements are in electronegativity.
    The aim is to give the transformer useful summaries
    that it would otherwise have to construct from the
    individual composition columns.

    \item \textbf{The original composition stays available.}
    The ten descriptors supplement the existing inputs:
    \[
    \begin{aligned}
        \text{Original inputs}&=17+3+1=21,\\
        \text{With descriptors}&=21+10=31.
    \end{aligned}
    \]
    The \(17\) element amounts remain separate columns.
    They are not replaced by an aggregate alloy score.

    \item \textbf{These are fixed calculations.}
    The same composition and lookup table always produce
    the same descriptors. Their calculation uses no EPIT
    labels and contains no separately learned material
    embedding or fitted corrosion model.
    \emph{Magpie-style} here refers to this specific
    ten-feature bundle.
\end{notes}

\subsection{From Weight Percentages to Atomic Fractions}
\label{sec:magpie-atomic-fractions}

\begin{notes}
    \item \textbf{Start from the raw composition.}
    Let \(c_{ij}\) be the weight percentage of element
    \(j\) in row \(i\), and let \(M_j\) be its atomic weight
    from the fixed lookup table.
    Descriptor calculation uses the raw element amounts,
    before feature standardization.

    \item \textbf{Calculate the fraction of atoms of each element.}
    Divide each weight percentage by its atomic weight,
    then normalize:
    \begin{equation}
        a_{ij}
        =
        \frac{c_{ij}/M_j}
             {\displaystyle\sum_{e=1}^{17}c_{ie}/M_e}.
        \label{eq:magpie-atomic-fractions}
    \end{equation}
    Here, \(a_{ij}\) is an atomic fraction, and
    \(\sum_{j=1}^{17}a_{ij}=1\).
    Equal weight percentages need not correspond to equal
    numbers of atoms because atomic weights differ.

    \item \textbf{The fractions describe the retained elements.}
    The sum includes only the \(17\) model-visible elements.
    Elements outside that subset do not contribute to the
    descriptors, even if they were present in the full
    source composition.
    The normalization therefore describes the retained
    composition subset.

    \item \textbf{The conversion uses a copy.}
    Atomic fractions are used only for descriptor calculation.
    The original composition columns retain their
    weight-percentage representation.
    The descriptor functions require finite, nonnegative
    amounts and a positive total amount in each row.
\end{notes}

\subsection{Descriptor Formulas and Their Meaning}
\label{sec:magpie-statistics}

\begin{notes}
    \item \textbf{Notation for an elemental property.}
    Let \(q_j\) be a tabulated property of element \(j\),
    such as electronegativity or melting temperature.
    Different descriptors apply the following statistics
    to different property tables.

    \item \textbf{Weighted mean: the average property per atom.}
    \begin{equation}
        \mu_q^{(i)}
        =
        \sum_{j=1}^{17}a_{ij}q_j.
        \label{eq:magpie-mean}
    \end{equation}
    Elements with larger atomic fractions contribute more.

    \item \textbf{Average absolute deviation: how spread out the properties are.}
    \begin{equation}
        \delta_q^{(i)}
        =
        \sum_{j=1}^{17}a_{ij}
        \left|q_j-\mu_q^{(i)}\right|.
        \label{eq:magpie-deviation}
    \end{equation}
    This measures the weighted average distance from the
    mean. It is an absolute deviation, not a variance or
    standard deviation.

    \item \textbf{Range: the difference between the largest and smallest properties.}
    \begin{equation}
        R_q^{(i)}
        =
        \max_{j:\,a_{ij}\geq0.01}q_j
        -
        \min_{j:\,a_{ij}\geq0.01}q_j.
        \label{eq:magpie-range}
    \end{equation}
    Only elements with at least \(1\%\) atomic fraction
    enter the range. This prevents a tiny trace of an
    element from determining an extreme.
    Means and deviations still include every element with
    a positive fraction; the cutoff does not change the
    composition itself.

    \item \textbf{A simple numerical example.}
    Suppose two elements have atomic fractions \(0.75\)
    and \(0.25\), and property values \(4\) and \(8\).
    Then
    \[
    \begin{aligned}
        \mu_q&=0.75(4)+0.25(8)=5,\\
        \delta_q&=0.75|4-5|+0.25|8-5|=1.5,\\
        R_q&=8-4=4.
    \end{aligned}
    \]
    These illustrative numbers show that the three
    statistics describe different aspects of the same
    composition.
\end{notes}

\subsection{The Ten Appended Features in Their Exact Order}
\label{sec:magpie-feature-order}

\begin{notes}
    \item \textbf{Columns \(22\)--\(24\): electronegativity.}
    Append its mean, range, and average absolute deviation,
    in that order.
    These summarize the tabulated electron-attracting
    tendencies of the constituent elements and their
    differences.

    \item \textbf{Columns \(25\)--\(26\): covalent radius.}
    Append its mean and average absolute deviation.
    These summarize the tabulated atomic sizes and their
    variation.

    \item \textbf{Columns \(27\)--\(28\): elemental melting temperature.}
    Append its mean and range.
    The mean is a weighted average of the elements'
    tabulated melting temperatures; it is not a prediction
    of the alloy's actual melting temperature.

    \item \textbf{Column \(29\): \(d\)-valence electrons.}
    Append the mean of the Magpie table's
    \(d\)-valence-electron counts.

    \item \textbf{Column \(30\): total valence electrons.}
    Append the mean of the table's total
    valence-electron counts.

    \item \textbf{Column \(31\): unfilled valence states.}
    Append the mean of the table's unfilled-valence-state
    counts. Like the other electron descriptors, this is
    a composition summary of fixed elemental values,
    not a calculation of the alloy's electronic structure.
\end{notes}

\subsection{Where the Descriptors Enter Synthetic Training}
\label{sec:magpie-synthetic-integration}

\begin{notes}
    \item \textbf{First construct the physical task and target.}
    Generate the \(21\) raw physical columns and calculate
    the SCM--EPIT response as described in
    Section~\ref{sec:informed-prior}.
    The SCM uses standardized physical inputs; the
    analytical rule uses their raw values.

    \item \textbf{Then append the composition summaries.}
    Calculate the ten descriptors from the raw composition:
    \[
        X_{\mathrm{raw}}^{31}
        =
        [X_{\mathrm{raw}}^{21},
         X_{\mathrm{Magpie}}^{10}].
    \]
    The superscripts indicate column counts.
    The target mechanisms still use only the original
    physical inputs, so adding the descriptors does not
    change the synthetic target for an otherwise identical
    generated task.

    \item \textbf{Standardize the full table column by column.}
    The transformer receives the standardized
    \(31\)-column table. With the procedure from
    Section~\ref{sec:feature-standardization},
    \[
        X_{\mathrm{transformer},1:21}^{31}
        =
        X_{\mathrm{std}}^{21}.
    \]
    Its first \(21\) columns therefore remain numerically
    identical to the standardized physical columns used
    by the SCM.

    \item \textbf{The descriptor columns have no separate random generator.}
    They receive no independent noise, SCM draw, shared
    block coupling, categorical conversion, or random
    marginal transformation.
    Their relationship with the target comes through the
    composition from which both are calculated.

    \item \textbf{Generic tasks receive no material descriptors.}
    The generic branch retains its own preprocessing.
    When mixed tasks need a common tensor width, unused
    positions can be filled with zeros for batching.
    Those padded positions are not active Magpie features.
\end{notes}

\subsection{Applying the Same Descriptors to Real Data}
\label{sec:magpie-real-data}

\begin{notes}
    \item \textbf{Prepare each evaluation split separately.}
    Fit mean imputation for the \(17\) composition columns
    using only the labeled context rows.
    Apply those fitted means to missing composition values
    in both context and query rows.

    \item \textbf{Use the same formulas and lookup values.}
    Convert the imputed weight percentages to atomic
    fractions and calculate the same ten descriptors
    in the order listed in
    Section~\ref{sec:magpie-feature-order}.
    Query EPIT labels are not needed for this calculation.

    \item \textbf{Keep the physical columns first.}
    Append the descriptors after the original \(21\)
    columns, then use the estimator's usual context-fitted
    preprocessing.
    Observed composition values are preserved; missing
    values have been replaced by the fitted means.
    The environment and method columns remain in their
    existing positions.

    \item \textbf{Match the model's trained representation.}
    A model trained with descriptors must receive the same
    augmentation during evaluation.
    The current \(21\)-input CorrPFN result is not an
    evaluation of this \(31\)-input extension.
\end{notes}

\subsection{Implementation, Validation, and Experimental Scope}
\label{sec:magpie-implementation}

\begin{notes}
    \item \textbf{Fixed source of elemental values.}
    The atomic weights and property values were extracted
    once from Matminer~0.10.0 and stored under the version
    identifier
    \path{epit_magpie_v1_matminer_0_10_0}.
    The calculation is implemented in
    \path{src/tabicl/prior/magpie_features.py}.
    It uses Torch during synthetic generation and NumPy
    during real-data evaluation; Matminer is not a runtime
    dependency.

    \item \textbf{How the optional configuration is recorded.}
    The training setting is
    \texttt{pitting\_magpie\_features}.
    The final-training workflow supports a recorded
    \texttt{--override-use-magpie} option.
    Such an override changes the effective training
    configuration; it does not mean that the original
    Optuna trial selected descriptors.

    \item \textbf{What the recorded checks cover.}
    Earlier validation in \texttt{main2.tex} reports
    agreement with Matminer reference values.
    The repository also contains checks for Torch/NumPy
    agreement, preserved base columns, context-fitted
    imputation, descriptor order, and \(31\)-feature
    informed generation.
    These checks concern correct calculation and integration.

    \item \textbf{Predictive benefit remains a separate comparison.}
    Correct descriptor values do not establish improved
    EPIT prediction. Historical descriptor experiments
    changed other parts of the generator and evaluation.
    Establishing the benefit for the final v7 construction
    requires a matched comparison with descriptors enabled
    and disabled, following the evaluation principles in
    Section~\ref{sec:next-experiments}.
\end{notes}

\end{document}
